\documentclass[11pt]{article}
\usepackage[margin=1in]{geometry}
\usepackage{amsmath,amssymb,amsthm,booktabs,longtable,array}
\usepackage[T1]{fontenc}
\usepackage{lmodern}
\usepackage{xcolor}
\usepackage{xurl}
\usepackage[colorlinks=true,linkcolor=blue!45!black,citecolor=blue!45!black,urlcolor=blue!45!black]{hyperref}
\hypersetup{pdftitle={A generating grammar for six-point homometry in finite cyclic groups},pdfauthor={Hunter Bown},pdfsubject={Computer-assisted six-point generation; review draft prepared with AI assistance}}
\usepackage{enumitem}
\newtheorem{theorem}{Theorem}[section]
\newtheorem{lemma}[theorem]{Lemma}
\newtheorem{proposition}[theorem]{Proposition}
\newtheorem{corollary}[theorem]{Corollary}
\newtheorem{definition}[theorem]{Definition}
\theoremstyle{remark}
\newtheorem{remark}[theorem]{Remark}
\newcommand{\Z}{\mathbb Z}
\newcommand{\R}{\mathbb R}
\newcommand{\C}{\mathbb C}
\newcommand{\Zn}{\Z/n\Z}
\newcommand{\TI}{\mathrm{T/I}}
\newcommand{\ord}{\operatorname{ord}}
\newcommand{\rank}{\operatorname{rank}}
\newcommand{\Hom}{\operatorname{Hom}}
\newcommand{\supp}{\operatorname{supp}}
\newcommand{\row}{\operatorname{row}}
\newcommand{\code}[1]{\nolinkurl{#1}}
\setlength{\emergencystretch}{3em}
\setlist{itemsep=3pt,topsep=5pt}
\title{A generating grammar for six-point homometry\protect\\in finite cyclic groups}
\author{Hunter Bown\\\small \texttt{hmbown@gmail.com}}
\date{4 October 2026\\\small Revision with named authorship and machine-checked lemmas}
\begin{document}
\maketitle
\begin{center}
\small\textbf{Review draft.} Computer-assisted proof prepared with substantial AI assistance;\newline
external specialist review and historical priority remain unresolved. Not submitted.
\end{center}
\begin{abstract}
Two subsets of a finite cyclic group are homometric when their unordered
interval multisets agree. We give an explicit generating grammar for
six-element subsets in every finite cyclic group: the connected components
of its graph on translation/inversion classes are exactly the homometry
classes. The grammar combines the classical two-parameter integer
six-point construction, conditional block translations and reflections, a
parallelogram-dyad exchange, half-per-coset complementation,
autocorrelation-preserving units and thirteen rigid cyclic templates.
The completeness argument uses universal signed edge-matchings. Their
finite presentations reduce to source orders at most 135; their positive
free-rank presentations are covered by a weighted real-line theorem and
finite integral quotient certificates. The proof specifies the matching,
height-arrangement and quotient algorithms and their acceptance conditions.
The computations have separately implemented checks, but remain explicit
trust dependencies. The largest
two-method six-subset census is through modulus 135. The result concerns
generation with compositions, rather than unique normal forms, minimal
generators or closed counting formulas. No novelty claim is made.
\end{abstract}
\tableofcontents

\section{The problem and the claim}
For a set of pitch classes, an interval inventory forgets where its intervals
occur. This loss of placement is the mathematical question studied here;
it does not assert equal harmonic function or identical recorded sound.
In twelve-tone equal temperament, C major $\{0,4,7\}$ and C minor
$\{0,3,7\}$ have the same interval classes $\{3,4,5\}$ and are connected
by the reflection $p\mapsto7-p$. They belong to one translation/inversion
class. The four-note sets $\{0,1,4,6\}$ and $\{0,1,3,7\}$ also have
equal interval inventories, but occupy different such classes.

The four-point classification is prior work, including Rosenblatt's theorem
joint with Berman \cite[Theorem 3.9]{rosenblatt}. Erickson and Jones give
seven five-point types, including a triple family \cite{ej}.
General constructions and finite enumeration also precede this work
\cite{jj,mandereau}. The exact two-parameter six-point formula and its
signed factorization appear as Family N in Yovanof's thesis
\cite[p.~156]{yovanof} and in Yovanof--Golomb \cite[p.~47]{yg}; the distinct-distance context of that paper is not used
to infer weighted completeness here. The repeated-distance example at
$(p,q)=(1,4)$ is attributed to Hosemann and Bagchi (1954) in
Lemke--Skiena--Smith \cite[Table~1]{lss}. Classical block translations
appear in the Callender--Hall handout \cite[Cases~3--4]{ch}, which credits
Bullough \cite{bullough1} for Cases~4a/4b; Soderberg's dual inversion
\cite{soderberg} and Patterson's complementation theorems and examples
\cite{patterson} are further antecedents of the block and complement moves. Goyette describes translation and reflection
formulations \cite[Formulas~3.2 and~3.5]{goyette}; his proposed criteria
are explicitly empirical, so they are not proof dependencies here.
The symbolic-matching treatment of the six-point integer mechanism is also
relevant \cite{pg}. Our objective is an explicit grammar
valid at arbitrary modulus, rather than a finite census or an unrestricted
spectral-unit existence statement.

The argument has two kinds of step. Mathematical lemmas prove soundness,
normalization, exhaustive reductions, specialization and path inflation.
Finite verification statements exhaust the resulting matching presentations
and certify explicit mechanisms. We specify the latter algorithms and
acceptance predicates as part of the proof, rather than treating reported
totals as evidence of exhaustiveness. Section~\ref{sec:certificates} identifies
the saved inputs, arithmetic witnesses and independent checker scopes.
The weighted real-line result also uses a precisely stated exact
unsatisfiability obligation. These computational dependencies are retained
in the final conclusion; a bounded census alone would not prove it.

The statement is a generation theorem. It does not by itself solve the
stronger classification-and-enumeration problem of choosing irredundant
parameters and deriving counts. We distinguish that remaining problem from
the completeness of the graph's generating grammar.

\section{Definitions and the generating graph}
Let $C_n=\Zn$, with $n$ a positive integer. For a finite subset $A\subset C_n$,
define its directed autocorrelation by
\[
 r_A(t)=\#\{(a,a')\in A^2:a-a'=t\}.
\]
For $1\le d\le\lfloor n/2\rfloor$, its interval-class count is the number
of unordered pairs at cyclic distance $d$. Away from an antipodal distance,
this count is $r_A(d)$; at $d=n/2$ it is $r_A(d)/2$.
For sets of the same cardinality, equality of interval inventories is
therefore equivalent to $r_A=r_B$. We call this relation homometry.

Translation and inversion act by $A\mapsto t+A$ and $A\mapsto t-A$.
Their orbits are $\TI$ classes. Two homometric sets in different $\TI$
classes are a \emph{Z-pair}. A nontrivial homometry family consists of at
least two such classes. Multiplication by a unit is \emph{not} included
in the definition of a vertex.

In the integral group ring $\Z[C_n]$, write a set for its indicator,
$x^s$ for translation by $s$, and $F^*$ for exponent reversal. Then
homometry is $AA^*=BB^*$. This group-ring identity concerns pitch-class
indicators; it is not an equality of audio spectra.

\begin{definition}[Admissibility]
Every endpoint produced by a construction below must contain exactly six
distinct elements. A named block partition must be disjoint. Independent
translations, inversions and interchange of endpoints are permitted.
When forming an edge, discard endpoints in the same $\TI$ class.
\end{definition}

Let $\mathcal G_n$ have the $\TI$ classes of six-subsets as vertices.
Join two distinct vertices if representatives satisfy one of the explicit
constructions in Section~\ref{sec:grammar}. Edges are undirected.

\begin{theorem}[Six-point generation; Theorem G]\label{thm:main}
For every positive integer $n$, two six-subsets of $C_n$ are homometric
if and only if their vertices in $\mathcal G_n$ are joined by a finite path.
Equivalently, the connected components are exactly the homometry classes.
\end{theorem}

The assertion is vacuous for $n<6$. Singleton components are allowed.
The theorem permits compositions; it does not assert one direct move for
every pair. It gives neither a disjoint parameter classification nor a
minimal generating set. Its in-house computer-assisted status and external
review boundary are specified in Section~\ref{sec:trust}.

\section{The explicit grammar}\label{sec:grammar}
\subsection{Four conditional block moves}
Choose a disjoint partition $A=U+W$ with $|U|+|W|=6$, and put $P=UW^*$.
The new moving block must again be disjoint from $U$.
\begin{center}
\small
\begin{tabular}{lll}
\toprule
Move & Endpoint $B$ & Required group-ring identity\\
\midrule
L2: periodic translation & $U+x^sW$ & $x^sP=P$\\
L3$^*$: cosymmetric translation & $U+x^sW$ & $x^{-s}P=P^*$\\
L4: halfturn translation & $U+x^sW$ & $2s=0$, $x^s(P+P^*)=P+P^*$\\
L5: block reflection & $U+x^cW^*$ & $UW^*=x^{-c}UW$\\
\bottomrule
\end{tabular}
\end{center}
These tests concern the small cross-difference distribution of the
chosen blocks. L3$^*$ does not assume that $s$ has a half in $C_n$;
replacing it by a halved-shift formula would lose legitimate cases.

\subsection{B: the classical six-point family}
For $p,q\in C_n$ define
\begin{align}
X&=\{0,p,q-2p,2q-2p,2q,3q-p\},\label{eq:bloomx}\\
Y&=\{0,p,q+2p,2q-p,2q+p,3q-p\}.\label{eq:bloomy}
\end{align}
Require admissibility and different $\TI$ classes. This formula and its signed factorization appear as Family N in
Yovanof--Golomb \cite[p.~47]{yg}. We retain the short name Bloom for this
mechanism; neither the formula nor the factor identity is claimed as new.
When both lists have six distinct integer points, integer $p,q$ give
line-homometric sets. Every admissible modular
image has an integer lift of that kind.

\subsection{D: a parallelogram-dyad exchange}
Choose $a,b\in C_n$ with four distinct corners
$R=\{0,a,b,a+b\}$, and let $\rho(t)=a+b-t$.
Choose nonnegative integer functions $U_0,V_0$ that are constant on
$\langle a\rangle$-cosets and $\langle b\rangle$-cosets, respectively,
with total combined mass twelve. Put $K=U_0+V_0$ and require
\begin{itemize}
\item $\rho(K)=K$ and $K=1$ on $R$;
\item $K(t)\in\{0,1,2\}$ off $R$;
\item $K(t)\in\{0,2\}$ at $\rho$-fixed points off $R$.
\end{itemize}
Choose a four-set $C$ off $R$ as follows. On a two-point $\rho$-orbit,
choose neither point for $K=0$, either point for $K=1$, and both for
$K=2$. At a fixed point choose it precisely when $K=2$. Set
\[
 X=C\cup\{0,a+b\},\qquad Y=C\cup\{a,b\}.
\]
The mass condition is
\[
12=\ord(a)\sum_{\langle a\rangle\text{-cosets}}U_0
 +\ord(b)\sum_{\langle b\rangle\text{-cosets}}V_0.
\]
At least one step has order at most twelve. Nonunit steps are permitted.
The periodic-weight decomposition is the reviewed cyclic kernel theorem D;
the construction does not ask an unrestricted homometry oracle for a partner.

\subsection{L7: half-per-coset complementation}
Choose a subgroup $H$ of even order $h\in\{2,4,6,12\}$, and $12/h$
distinct $H$-cosets. In each select $h/2$ points. Let $A$ be their union,
let $S$ be the union of the full selected cosets, and set $B=S\setminus A$.
The $n=h=12$ case is classical hexachord complementation.

\subsection{L6: a conditional unit move}
For a unit $u$ modulo $n$, set $B=uA$ only when
\[
 r_A(ut)=r_A(t)\qquad\text{for every }t\in C_n.
\]
The condition is essential: an arbitrary unit need not preserve the
interval vector. It is an actual edge between $\TI$ vertices, rather than
a silent change to their equivalence relation. L6 is used only in the
bounded finite branch of the completeness proof.

\subsection{R: thirteen rigid cyclic templates}\label{sec:rigid}
For a row $(q,X,Y)$ of Table~\ref{tab:rigid}, choose $t\in C_n$ with
$qt=0$, and produce $tX,tY$. The allowed orders are precisely those giving
six-distinct and $\TI$-distinct endpoints. Subgroup inflation is included.
\begin{table}[htbp]
\centering\small
\begin{tabular}{rlll}
\toprule
$q$ & $X$ & $Y$ & $\ord(t)$\\
\midrule
17 & $0,1,2,3,8,12$ & $0,1,2,6,7,9$ & 17\\
19 & $0,1,2,3,6,10$ & $0,1,2,4,5,11$ & 19\\
21 & $0,1,2,4,7,14$ & $0,1,3,7,8,10$ & 21\\
21 & $0,1,2,5,6,15$ & $0,1,2,6,7,10$ & 21\\
21 & $0,1,3,7,10,15$ & $0,1,4,7,14,16$ & 21\\
23 & $0,1,2,3,7,17$ & $0,1,2,4,17,18$ & 23\\
24 & $0,1,2,5,7,16$ & $0,1,2,6,9,11$ & 12,24\\
27 & $0,1,2,3,7,19$ & $0,1,2,3,8,12$ & 27\\
27 & $0,1,2,6,19,22$ & $0,1,3,17,21,22$ & 27\\
28 & $0,1,2,4,12,23$ & $0,1,3,5,11,12$ & 14,28\\
30 & $0,1,2,6,19,22$ & $0,1,3,9,13,14$ & 15,30\\
30 & $0,1,3,5,12,25$ & $0,1,6,9,11,13$ & 15,30\\
31 & $0,1,2,5,11,19$ & $0,1,2,6,20,23$ & 31\\
\bottomrule
\end{tabular}
\caption{A sufficient rigid-template list; neither minimality nor independence
from the other moves is asserted.}\label{tab:rigid}
\end{table}

\section{Soundness and specialization}
\begin{proposition}\label{prop:sound}
Every admissible grammar edge preserves autocorrelation.
\end{proposition}
\begin{proof}
For a block partition,
\[
 AA^*=UU^*+WW^*+P+P^*.
\]
Translation of $W$ replaces its cross sum by $x^{-s}P+x^sP^*$ while
preserving $WW^*$. L2 leaves both cross terms unchanged; L3$^*$ exchanges
them; L4 preserves their sum because $s=-s$. For L5 the first new cross
term is $x^{-c}UW=P$, and its involution supplies the second.

For B use the Laurent polynomials
\[
 F=1+u+v,\quad Q=u^2+uv^2-uv+v,\quad
 R=u^2v-uv+u+v^2=u^2v^2Q^*.
\]
Direct expansion gives $X=u^{-2}FQ$, $Y=u^{-1}FR$ under
$u=x^p,v=x^q$. Therefore $XX^*=YY^*$. The signed factor and its
cancellations are essential. A collision-free $0/1$ direct-sum flip
with factor cardinalities two and three is rigid: every two-point factor
is a translate of its reflection, so flipping either factor gives only
a translation or reflection of the whole set.

For D, the selections have
$K=C+x^{a+b}C^*+(1+x^a)(1+x^b)$. Its mass is $2|C|+4=12$,
so $|C|=4$. Expansion gives
\begin{equation}\label{eq:dyad}
 XX^*-YY^*=x^{-a-b}(1-x^a)(1-x^b)K.
\end{equation}
The two periodic weights are annihilated by the displayed product.
For L7, the half-coset condition gives $SA^*+AS^*=SS^*$, hence
$(S-A)(S-A)^*=AA^*$. L6 permutes the prescribed equal coefficients.
For R, counting the fifteen unsigned differences in each table row
gives equality in $C_q$, which survives $x\mapsto t$.
\end{proof}

Homomorphic transport of homometry has classical antecedents
\cite{mandereaua}. All group-ring identities transfer through homomorphisms. To use a
transferred identity as a named generator, its combinatorial conditions
must also survive. Injectivity on each endpoint preserves block disjointness.
A halfturn that becomes zero gives an equal pair. For D, if the images of
the equal-sum dyads share one corner, they share both, so the two endpoints
coincide. Every nontrivial image retains four distinct corners, and the
cyclic periodic-kernel result D supplies its periodic weights.

\begin{lemma}[The L7 degeneration lemma]\label{lem:l7}
Under a homomorphism injective on each endpoint of a six-point L7 pair,
every $\TI$-distinct image is itself literal L7.
\end{lemma}
\begin{proof}
If the kernel on $H$ has size $k$, a half-coset's $h/2$ points must
inject into $h/k$ positions, so $k\le2$. For $k=2$, both complementary
halves fill the image coset; the two entire endpoints agree. Thus $k=1$
in a nontrivial image.

At most two occupied source cosets can now merge, since each contributes
$h/2$ distinct points to an image coset of size $h$. A merger gives one
full common image coset. For $h=2$, translation by the nonzero element
of $H$ exchanges every remaining singleton half and fixes every full
coset; the whole pair is translation equivalent. For $h=4$ there are
three occupied cosets. A merger leaves one complementary dyad; each
dyad in $C_4$ translates to its complement, with the same translation
fixing the common full coset. For $h=6$, merging exhausts the two cosets
and makes the endpoints equal. For $h=12$ there is only one coset.
Thus a $\TI$-distinct image has no merger and retains literal L7.
\end{proof}

\section{The dyad kernel and complete shape parameterization}\label{sec:dyadfull}

The periodic weights in D can be generated without solving an inverse
homometry problem. We give the kernel theorem and the recovery argument
explicitly. Throughout this section, functions have integer coefficients
on $C_n$; translation acts by $x^aW(t)=W(t-a)$.

\begin{lemma}[Nonnegative periodic decomposition]\label{lem:periodic}
Let $W\colon C_n\to\Z_{\ge0}$ and $a,b\in C_n$. Then
\[
 (1-x^a)(1-x^b)W=0
 \quad\Longleftrightarrow\quad
 W=U+V,
\]
where $U,V$ are nonnegative integer functions satisfying
$x^aU=U$ and $x^bV=V$. If $W$ has total mass $m$, then
\[
 m=\ord(a)\sum_{C_n/\langle a\rangle}U
   +\ord(b)\sum_{C_n/\langle b\rangle}V,
\]
where the sums take one value per coset.
\end{lemma}
\begin{proof}
The reverse direction follows by annihilating the two summands. For the
forward direction use the characters $\chi_k(t)=\exp(2\pi i kt/n)$,
$0\le k<n$. Translation is diagonal on these characters, and a product
eigenvalue $(1-\chi_k(-a))(1-\chi_k(-b))$ vanishes precisely when at least
one of its factors does. Hence the complex kernel of the product is the
sum of the two periodic subspaces. Taking real parts gives real periodic
functions $U_0,V_0$ with $W=U_0+V_0$.

This alone does not give nonnegative integral summands. Put
$H=\langle a\rangle$, $J=\langle b\rangle$ and $L=H+J$, and work
in one coset of $L$. List its $H$-cosets by $i$ and its $J$-cosets by $j$.
Every $H$-coset meets every $J$-coset: the difference between their
representatives belongs to $H+J$. All intersections are nonempty, have
the same size $|H\cap J|$, and partition the component. The function
$W=U_0+V_0$ is constant on an intersection; call its integer value $w_{ij}$.
It satisfies the rectangle identities
\[
 w_{ij}-w_{ij_0}=w_{i_0j}-w_{i_0j_0}
 \qquad\text{for all }i,j.
\]
Choose $j_0$ and set $m_0=\min_i w_{ij_0}$,
\[
 u_i=w_{ij_0}-m_0,\qquad
 v_j=w_{i_0j}-w_{i_0j_0}+m_0.
\]
These numbers are integral, $u_i\ge0$, and the rectangle identities give
$v_j=\min_i w_{ij}\ge0$ and $w_{ij}=u_i+v_j$. Assign $u_i$ to its whole
$H$-coset and $v_j$ to its whole $J$-coset. Repeat on each $L$-coset.
The resulting $U,V$ are the required nonnegative integral functions.
Summing their constant values along cosets proves the mass formula.
\end{proof}

\begin{proposition}[Exactness of D within the dyad shape]\label{prop:dyadshape}
Suppose independent rigid motions align a six-pair to
\[
 X=C\sqcup\{0,a+b\},\qquad Y=C\sqcup\{a,b\},
\]
where $|C|=4$ and $R=\{0,a,b,a+b\}$ consists of four distinct elements
outside $C$. The pair is homometric if and only if it has the periodic
weights and orbit selections specified in construction D.
\end{proposition}
\begin{proof}
Write $\rho(C)=x^{a+b}C^*$ and $K=C+\rho(C)+1_R$.
The expansion \eqref{eq:dyad} and invertibility of $x^{-a-b}$ show that
homometry is equivalent to $(1-x^a)(1-x^b)K=0$. The function $K$ is
nonnegative, integral, has mass twelve, and is fixed by $\rho$.
Its coefficients are one on $R$ and belong to $\{0,1,2\}$ off $R$.
At a $\rho$-fixed point off $R$ its coefficient is twice the coefficient
of $C$, and therefore belongs to $\{0,2\}$. Lemma~\ref{lem:periodic}
supplies exactly the required weights.

Conversely, admissible weights give a $K$ with these properties. On each
two-point $\rho$-orbit, choose no point, one point, or both points according
as its coefficient is zero, one, or two. At a fixed point choose it
precisely when its coefficient is two. This produces a binary set $C$
off $R$ with $C+\rho(C)=K-1_R$. Summing coefficients gives
$2|C|=12-4=8$, so its cardinality is four. The dyad identity now gives
homometry. The same selection rule recovers every $C$ in the forward
direction. The final $\TI$ test decides whether it is a Z-pair.
\end{proof}

If $\ord(a)>12$ then its periodic summand must vanish, because a nonzero
nonnegative integer weight would alone contribute more than twelve.
Thus $K$ is $b$-periodic and $\ord(b)$ divides twelve. If both orders
exceed twelve, the construction is impossible. These are consequences
of the exact mass formula, not search caps. This proposition classifies
the displayed exchange shape; it does not assert that every six-pair
can be aligned to that shape.

\subsection{Coefficient proof of half-per-coset complementation}
Let $H$ have even order $h$, and let $A$ contain $h/2$ points in each
of $m=12/h$ selected $H$-cosets. Write $S$ for their full union. For
any $t\in C_n$, let $N_t$ count the ordered pairs of selected cosets
$(K,J)$ with $K=t+J$. Then direct coefficient counting gives
\[
 (SS^*)(t)=hN_t,\qquad
 (AS^*)(t)=(SA^*)(t)=\tfrac h2N_t.
\]
For example, for each such pair $(K,J)$ exactly $h/2$ choices of a point
of $A\cap K$ have their difference-$t$ partner in $J$. Consequently
$AS^*+SA^*=SS^*$, and expansion of $(S-A)(S-A)^*$ proves L7.
The equality of the two cross coefficients here follows from the
half-coset hypothesis, rather than from a general symmetry assumption
on $A$. The number of selected cosets is integral precisely for
$h\in\{2,4,6,12\}$ in the six-point construction.


\section{Universal signed matchings and the original presentation}
\label{sec:universal}

We distinguish a relation in an integral presentation from a relation that
holds only after a particular specialization. This distinction is essential:
dividing an integral relation by its content can remove precisely the torsion
that creates a cyclic Z-pair.

\begin{lemma}[Occurrence matching]\label{lem:occurrence-matching}
Let $(a_i)_{i=0}^5$ and $(b_i)_{i=0}^5$ be two lists in an abelian group,
including multiplicities. Their directed autocorrelations agree if and only
if there are a bijection $\sigma$ of their fifteen unordered edge
occurrences and signs $\epsilon_{ij}\in\{1,-1\}$ such that
\[
 a_j-a_i=\epsilon_{ij}(b_l-b_k),\qquad
 (k,l)=\sigma(i,j),\quad i<j, k<l.
\]
\end{lemma}
\begin{proof}
Each unordered edge with difference $g$ contributes the two directed
occurrences $g,-g$. Partition differences into orbits under $g\mapsto-g$.
For an orbit of size two, its autocorrelation coefficient at either element
is the number of unordered edges in that orbit. For a nonzero orbit of size
one it is twice that number. At zero it is six plus twice the number of
coincident off-diagonal unordered edges. Consequently equality of
autocorrelations gives equal numbers of edge occurrences in every orbit.
Match those occurrences bijectively within each orbit and choose a sign
that gives equality. If $g=-g$, both signs are available and neither may be
discarded. Conversely, a signed bijection equates the two directed
occurrences contributed by each edge, and the six diagonal zeros agree.
\end{proof}

Write $\mathcal E=\{(i,j):0\le i<j\le5\}$. In $\Z^{10}$ use the row vectors
$e_1,\ldots,e_5,f_1,\ldots,f_5$ as the standard basis, and set
$e_0=f_0=0$. Independently translate each input endpoint so that
$a_0=b_0=0$. A signed matching defines the $15\times10$ matrix with rows
\begin{equation}\label{eq:formal-matching-row}
 m_{ij}=e_j-e_i-\epsilon_{ij}(f_l-f_k),\qquad(k,l)=\sigma(i,j).
\end{equation}
The \emph{integral relation lattice}, the \emph{original presentation},
and its free rank are
\begin{equation}\label{eq:original-presentation}
 L_M=\sum_{(i,j)\in\mathcal E}\Z m_{ij},\qquad
 G_M=\Z^{10}/L_M,\qquad d=10-\rank_{\mathbb Q}M.
\end{equation}
We write $E_i,F_i$ for the classes of $e_i,f_i$ in $G_M$. The two labelled
lists $(E_i)$ and $(F_i)$ are called the universal configurations.

\begin{lemma}[Universal property and rank]\label{lem:universal-property}
The universal configurations are homometric as lists. Every realization
of the signed matching in an abelian group $K$ factors through a
homomorphism $G_M\to K$. Moreover $5\le\rank_{\mathbb Q}M\le10$, so
$0\le d\le5$. A realization by six-element sets rules out universal
within-endpoint collisions; a $\TI$-distinct realization rules out
universal $\TI$ equivalence.
\end{lemma}
\begin{proof}
Equation~\eqref{eq:formal-matching-row} gives a signed edge bijection in
$G_M$, so Lemma~\ref{lem:occurrence-matching} gives homometry. Sending
$e_i$ to $a_i$ and $f_i$ to $b_i$ defines a map from $\Z^{10}$ that kills
every row of $M$, hence factors through the quotient. The first five
columns of $M$ are the reduced incidence matrix of the complete graph
on six vertices; the rows for edges $(0,i)$ contain the five standard
basis vectors, giving rank five. There are ten columns. Finally a
universal collision or rigid equivalence remains true under every
homomorphism and therefore under the given realization.
\end{proof}

The rank in~\eqref{eq:original-presentation} belongs to the original
signed matching. A map from $G_M$ into a cylinder or a finite cyclic group
can impose additional relations and lower its rank; this never changes
which branch the original presentation occupies. In particular, we never
seek a nonzero map from the finite target $C_n$ to $\R$.

\subsection{The exact meaning of a lattice certificate}
For row matrices $M,N$ with ten columns, two integer matrices satisfying
\begin{equation}\label{eq:two-way-lattice-cert}
 AM=N,\qquad BN=M
\end{equation}
certify $L_M=L_N$. An equality of rational row spaces does not certify this
claim. The latter equality would identify only the free solutions and
could identify different torsion groups. We use primitive rational rows
only when enumerating height subspaces below.

To certify coordinates for $G_M$, supply integer matrices $U,V,D$ with
\begin{equation}\label{eq:smith-contract}
 UMV=D,\qquad \det U,\det V\in\{1,-1\},
\end{equation}
where $U$ is $15\times15$, $V$ is $10\times10$, and $D$ is rectangular
diagonal with nonzero entries $d_1,\ldots,d_r$ and all other entries zero.
The map sending a row vector $z$ to $zV$ induces an isomorphism
\[
 G_M\cong\bigoplus_{i=1}^{r}\Z/|d_i|\Z\ \oplus\ \Z^{10-r}.
\]
Factors of order one are omitted. Thus the coordinates of an old generator
are its row of $V$, reduced modulo $|d_i|$ in each torsion position. All
of~\eqref{eq:two-way-lattice-cert}--\eqref{eq:smith-contract} can be
checked by integer multiplication and determinants. Discovery by Hermite
or Smith algorithms is useful; their claimed normal forms are not accepted
without these identities. A collision certificate can instead be a row
$c$ with $cM=e_i-e_j$ (or $f_i-f_j$), $i\ne j$. A rigid-equivalence
certificate consists of the corresponding six integer row identities.

\subsection{The finite torsion bound and the finite branch}
\begin{lemma}\label{lem:bound}
The torsion subgroup $T_M$ of the original presentation $G_M$ has order
at most $135$.
\end{lemma}
\begin{proof}
First every square minor of a reduced oriented incidence matrix has
determinant $0,1$ or $-1$. In a chosen square submatrix each row has
at most two entries; if some row has at most one nonzero entry, expansion
reduces the assertion by induction. Otherwise every row has two opposite
entries, so the sum of its columns is zero and its determinant vanishes.
Permuting or changing signs of rows preserves this property. Both
five-column blocks of $M$ have it.

An $r\times r$ minor using $c$ columns of the first block has a Laplace
expansion with $\binom rc$ summands, each a product of two incidence
minors of absolute value at most one. For $r\le9$ this is at most
$\binom94=126$. If $r=10$, all ten columns occur. The chosen ten rows
correspond to ten distinct edges of the first $K_6$. A five-row
first-block determinant is nonzero precisely for a spanning tree:
a cycle gives dependence; five acyclic edges on six vertices form a
connected tree; and leaf elimination gives determinant $\pm1$ for a
tree. The complementary second-block determinant still has absolute
value at most one. Therefore
\[
 |\det M_R|\le\tau(K_6[R]),\qquad |R|=10,
\]
where $\tau$ counts spanning trees. No assertion that these summands
have a common sign is needed.

The bound on $\tau$ is an explicit small finite obligation with no
symmetry reduction. Enumerate every $R\subset\mathcal E$ of size ten
and every $J\subset\mathcal E$ of size five. Check connectedness and
acyclicity of the graph with edges $J$, and for every tree count its
containment $J\subset R$. Independently, form the $6\times6$ graph
Laplacian of $R$, remove one fixed row and column, and compute the exact
$5\times5$ determinant. The two answers agree because, for the reduced
incidence matrix $B_R$, Cauchy--Binet gives
\[
 \det(B_R^{\mathsf T}B_R)
   =\sum_{\substack{J\subset R\\|J|=5}}(\det B_J)^2
   =\tau(K_6[R]).
\]
This also proves the Laplacian method's counting interpretation, rather
than using an unproved counting shortcut.

The acceptance predicate requires every one of the $\binom{15}{10}=3003$
labelled $R$ and every one of the $\binom{15}{5}=3003$ possible $J$,
and equality of the two counts for each $R$. The complete histogram is:
\begin{center}\small
\begin{tabular}{r|rrrrrrr}
\toprule
$\tau$ & $0$ & $75$ & $96$ & $99$ & $100$ & $104$ & $111$\\
Number of $R$ & $6$ & $300$ & $90$ & $360$ & $90$ & $360$ & $360$\\
\midrule
$\tau$ & $114$ & $115$ & $120$ & $121$ & $128$ & $130$ & $135$\\
Number of $R$ & $360$ & $360$ & $180$ & $72$ & $45$ & $360$ & $60$\\
\bottomrule
\end{tabular}
\end{center}
The counts sum to $3003$, and the maximum is $135$. The project and
the independent review reproduce the full histogram by the two methods.
The incidence argument together with this exact finite obligation bounds
every nonzero maximal-rank minor by $135$.

Finally the ideal generated by the rank-$r$ minors is unchanged by
unimodular row and column operations, by Cauchy--Binet applied to those
operations and their integer inverses. In a diagonal presentation it
is generated by $\prod_{i=1}^r d_i$, whose absolute value is $|T_M|$.
Thus $|T_M|$ is the positive gcd of the nonzero maximal-rank minors.
It is no greater than any one of their absolute values and therefore
at most $135$.
\end{proof}

\begin{corollary}\label{cor:finite}
If $d=0$ and $G_M\to C_n$ realizes six-element endpoints, then, after
their independent anchors are removed, both endpoints lie in a cyclic
subgroup of order $q$ with $6\le q\le135$ and $q\mid n$. They are an
ordinary subgroup inflation of two homometric six-subsets of $C_q$.
\end{corollary}
\begin{proof}
Now $G_M=T_M$ is finite. Its image $H$ in $C_n$ is a cyclic subgroup,
so its order $q$ divides $n$ and $|T_M|$, and hence is at most $135$.
Six distinct endpoint coordinates force $q\ge6$. Identify $H$ with
$C_q$ by the injective map $t\mapsto(n/q)t$. The inverse on $H$
transports the signed matching and the six-element endpoints to $C_q$.
Restoring the separate translations of the original endpoints does not
change their vertices in the grammar graph.
\end{proof}

\subsection{The real solution space and a fixed congruence}
Let
\[
 V_M=\Hom(G_M,\R)=\{v\in\R^{10}:Mv=0\}.
\]
This real vector space has dimension $d$. Its points are the coordinates
of real homometric six-atom lists, even when distinct universal vertices
have identical real coordinates.

\begin{lemma}[Finite unions of subspaces]\label{lem:finite-subspaces}
Over an infinite field a finite-dimensional vector space cannot be covered
by finitely many proper linear subspaces.
\end{lemma}
\begin{proof}
For each proposed proper subspace choose a nonzero linear functional
vanishing on it. Their product is a nonzero polynomial in a basis of the
whole space. A nonzero polynomial cannot vanish at all points over an
infinite field: in one variable it has only finitely many roots; induction
on the number of variables applies this fact to its coefficient
polynomials. The product therefore cannot vanish throughout the space.
\end{proof}

For a permutation $\pi$ of six labels and $\eta\in\{1,-1\}$, define
\[
 T_{\pi,\eta}=\{v\in\R^{10}:
  b_i=\eta(a_{\pi(i)}-a_{\pi(0)})\text{ for every }i\}.
\]
These are linear subspaces. Since $a_0=b_0=0$, every real translation or
reflection equivalence of the two multisets has exactly this form for
some $\pi,\eta$. Labels with equal coordinates can be matched arbitrarily;
there are still only $2\cdot6!$ choices.

\begin{lemma}[Fixed free-quotient alignment]\label{lem:fixed-alignment}
If the two real multisets are congruent for every $v\in V_M$, then one
fixed $T_{\pi,\eta}$ contains $V_M$. After relabelling, reflecting and
reanchoring the second universal configuration, every real solution
satisfies $b_i=a_i$ for all $i$. These changes are invertible integral
changes of generators and preserve $G_M$ and $d$.
\end{lemma}
\begin{proof}
The hypothesis says that the finitely many subspaces
$V_M\cap T_{\pi,\eta}$ cover $V_M$. Lemma~\ref{lem:finite-subspaces}
implies that one is the entire space. Relabelling and reflection give
$b'_i=\eta(b_{\pi^{-1}(i)}-b_{\pi^{-1}(0)})=a_i$ on $V_M$.
A permutation followed by reanchoring transforms the five coordinates
of a labelled endpoint by an integer linear map. Undoing the permutation
and reanchoring is its integer inverse. Reflection multiplies that block
by $-1$. Hence the full ten-coordinate transformation is unimodular.
\end{proof}

The conclusion identifies the free projections, not the group elements
$E_i,F_i$. In a finitely generated abelian group an element is killed by
every real homomorphism precisely when it is torsion: write the group as
$\Z^d\oplus T$ and use its $d$ coordinate maps. Consequently the aligned
$F_i-E_i$ may be nonzero torsion. Imposing $F_i=E_i$ as \emph{integral}
relations would discard the cyclic phenomena under investigation.

Substitute $b_i=a_i$ only in the real equations. With six unanchored
heights, the resulting rows have the form
\begin{equation}\label{eq:height-rows}
 (u_j-u_i)-\epsilon(u_l-u_k),\qquad i<j, k<l,
\end{equation}
where $u_0,\ldots,u_5$ denote the six standard vectors. They have coordinate
sum zero. Let $S$ be their rational span. Removing the common translation
by setting $a_0=0$ gives a five-dimensional height space, and
\begin{equation}\label{eq:height-codimension}
 \dim_{\mathbb Q}S=5-d.
\end{equation}
Indeed every vector in $V_M$ lies on the aligned graph. Conversely any
anchored height vector satisfying~\eqref{eq:height-rows}, placed on that
graph, satisfies the original real matching equations and lies in $V_M$.
Thus the anchored kernel of $S$ is exactly $V_M$.

\section{The weighted real-line input and the Bloom branch}
\label{sec:formal-bloom}
This section states explicitly the mathematical formula checked by the
line solvers. The result is the previously reviewed I/Iw component, tagged
\textbf{[PROVED], in-house, computer-assisted}. Its completeness remains
conditional on the saved exact UNSAT results and their audited encodings;
the elementary interface proof below does not remove that dependency.

\begin{lemma}[Weighted six-atom line theorem Iw]\label{lem:iw}
Two noncongruent homometric real multisets of total multiplicity six
have, after independent translations and reflections and possible
interchange, one of the following weakly ordered forms:
\begin{align}
 A_1&=(0,q-2p,q,2q-3p,3q-2p,3q-p),\notag\\
 B_1&=(0,q-p,q+p,2q+p,3q-2p,3q-p),
       &&p>0,\quad q\ge3p;\label{eq:weighted-form-one}\\
 A_2&=(0,q-p,q+p,3q-2p,3q-p,2q+p),\notag\\
 B_2&=(0,p,2q-p,q+2p,3q-p,2q+p),
       &&p>0,\quad 3p/2\le q<2p.\label{eq:weighted-form-two}
\end{align}
For six-element sets the two lower bounds are strict. For integer
coordinates the recovered $p,q$ are integers. Every displayed pair is
homometric and noncongruent.
\end{lemma}

\begin{proof}[Mathematical and computational dependency]
The normalization, exact occurrence-count encoding, coefficient-plane
exclusions and solver obligation are given in Section~\ref{sec:linefull}.
The saved exact UNSAT results exclude every normalized pair outside the
two displayed planes over the unbounded real domain. Their weak-order
inequalities give precisely the parameter chambers above; the strict-order
version gives the strict lower bounds. Integer parameters are recovered
from $p=A_5-A_4,q=A_2$ in the first form and
$p=A_4-A_3,q=A_1+p$ in the second. The corresponding gap lists also prove
noncongruence. This is the reviewed Iw dependency, with its exact solver
trust retained, rather than a new solver run or a bounded-data inference.
\end{proof}

\subsection{The exact weighted real-line obligation}\label{sec:linefull}

The following exposition states the complete mathematical formula whose
unsatisfiability is used in BF. It is included so that the real-domain
claim can be evaluated independently of a program's normal-form names.
A six-atom multiset is a list of six real positions with repetitions
allowed; the fifteen unordered pairs refer to distinct labelled atoms,
even when their positions coincide.

\begin{lemma}[Diameter normalization]\label{lem:linenorm}
Every noncongruent homometric pair of six-atom real multisets has positive
common diameter $L$. After independent translations, reflections and
division by $L$, followed if necessary by interchange of endpoints, it
has the form
\[
 A=(0,a_1,a_2,a_3,r,1),\qquad B=(0,b_1,b_2,b_3,r,1),
\]
with weakly increasing coordinates, $r\ge1-a_1$, $r\ge1-b_1$, and
$(a_1,a_2,a_3)<_{\rm lex}(b_1,b_2,b_3)$. The two lists are not related
by reflection about $1/2$.
\end{lemma}
\begin{proof}
The largest labelled pair distance is the diameter, so homometry gives
the same diameter for both lists. Diameter zero would give two congruent
lists with all six atoms at one point. Anchor the minima at zero and
the maxima at one. Remove one labelled endpoint-edge occurrence of
length one from the common multiset of distances. In either list the
largest remaining value is $\max(a_4,1-a_1)$: both those endpoint
distances remain, and every other edge is no larger than one of them.
This argument still holds if an endpoint is repeated or the remaining
maximum is also one. Reflect each list independently to put its shared
remaining maximum at coordinate four. The two lists are unequal, since
equality would be a congruence; lexicographic interchange is therefore
available. Reflection is excluded by the original noncongruence.
\end{proof}

Put $z=(a_1,a_2,a_3,r,1,b_1,b_2,b_3,r,1)$. The counterexample formula
$\Phi(z)$ is the conjunction of the following predicates:
\begin{enumerate}
\item Each displayed list is weakly increasing and the two inequalities
$r\ge1-a_1$, $r\ge1-b_1$ hold.
\item The strict lexicographic inequality of Lemma~\ref{lem:linenorm}
holds, and
$\neg\bigwedge_{i=1}^{4}(A_i=1-B_{5-i})$ holds.
\item With $d_{ij}=A_j-A_i$ and $e_{kl}=B_l-B_k$, for every labelled
$i<j$,
\[
 \sum_{k<l}[d_{kl}=d_{ij}]
   =\sum_{k<l}[e_{kl}=d_{ij}].
\]
Here $[P]$ is one if $P$ holds and zero otherwise.
\item The vector $z$ belongs to neither of the two linear planes given
by the two height lists displayed in BF.
\end{enumerate}
The indices in the last condition run through the ten coordinates after
omitting the two anchored zeros. If a template has coefficient pairs
$(u_k,v_k)$, choose $i,j$ with $\Delta=u_iv_j-v_iu_j\ne0$. Membership
in its plane is the conjunction, for all ten indices $k$, of
\[
 \Delta z_k=(u_kv_j-v_ku_j)z_i+(-u_kv_i+v_ku_i)z_j.
\]
Thus nonmembership is the negation of a specified finite conjunction;
there is no bounded parameter grid or rationality condition.

\begin{lemma}[Encoding equivalence]\label{lem:lineencode}
$\Phi$ is satisfiable over the reals if and only if there is a normalized
noncongruent homometric six-atom pair outside both displayed planes.
\end{lemma}
\begin{proof}
Only the distance-count and plane predicates need explanation. The
distance equations give equality of multiplicity for every value occurring
among the fifteen $A$ edges, including zero. The required multiplicities
sum to fifteen, so there is no room for a further $B$ value. Conversely,
equal multisets satisfy every equation. Equalities among labelled atoms
are retained; the directed autocorrelation adds six diagonal zeros and
twice the unordered zero distances to both sides. Cramer's rule shows
that the displayed plane equations are exactly the existence of real
parameters $p,q$ with $z_k=u_kp+v_kq$. The other predicates are the
normalization and exclusion conditions already proved.
\end{proof}

\paragraph{Finite solver obligation U.}
The exact formula $\Phi$ is unsatisfiable. This is the saved exact
linear-real-arithmetic obligation checked by the weighted builder and
a separate consecutive-gap/occurrence-bijection encoding. Its truth at
all real positions is a solver dependency of this computer-assisted
proof. A bounded integer recount does not establish U. Saved Z3 and
cvc5 results, input hashes and internal proof-checking settings are
specified in the evidence inventory; their proof objects have not been
translated into an independent general-purpose proof assistant.

Accepting U, Lemmas~\ref{lem:linenorm} and~\ref{lem:lineencode} put every
pair in one of the two planes. Recovering the parameters in the first
plane gives $p=A_5-A_4$, $q=A_2$. The weak consecutive gaps force
$p\ge0$ and $q\ge3p$; $p=0$ makes the lists equal. In the second plane,
$p=A_4-A_3$, $q=A_1+p$; the gaps force $p>0$,
$3p/2\le q\le2p$, and the upper endpoint makes the lists equal.
Thus the exact domains are the relaxed chambers stated in BF. Their
formal positive-distance identities prove soundness even on the allowed
repeated-coordinate boundaries. For completeness their consecutive gap
lists are, respectively,
\begin{align*}
 A_1 &: (q-2p,2p,q-3p,q+p,p),\\
 B_1 &: (q-p,2p,q,q-3p,p),\\
 A_2 &: (q-p,2p,2q-3p,p,2p-q),\\
 B_2 &: (p,2q-2p,3p-q,2q-3p,2p-q).
\end{align*}
The first coordinates after zero in the first pair differ by $p>0$.
Reflecting $A_1$ makes that coordinate $p$, whereas $B_1$ has
$q-p\ge2p$. In the second pair the two coordinates are $q-p$ and $p$,
which differ because $q<2p$; the reflected $A_2$ has coordinate
$2p-q\le p/2<p$. Anchored equal-diameter lists can be congruent only by
equality or this reflection. Thus both forms are noncongruent, including
their allowed repeated-coordinate boundaries. Integer input coordinates recover integer
$p,q$ after undoing the diameter normalization: the formulas use only
integer point differences. For distinct six-sets the lower bounds become
strict. This is the full logical passage from U to Iw and I.


\begin{lemma}[BF: noncongruent real projection]\label{lem:bf}
Let $A,B$ be homometric six-element subsets of an abelian group $K$.
If some $\ell:K\to\R$ sends their six-atom lists to noncongruent
multisets, then, up to independent rigid motions and interchange,
$A,B$ have the group-valued Bloom formula
\[
 X=\{0,p,q-2p,2q-2p,2q,3q-p\},\qquad
 Y=\{0,p,q+2p,2q-p,2q+p,3q-p\},\quad p,q\in K.
\]
The conclusion depends on Iw and the finite integral certificate predicate
described in this proof.
\end{lemma}
\begin{proof}
Use Iw to normalize the two projected lists. The needed translations
subtract actual elements of $A$ or $B$, and reflections use inversion
and such translations, so they can be performed in $K$. Label equal
height atoms arbitrarily. The resulting projected coordinates have one
of~\eqref{eq:weighted-form-one}--\eqref{eq:weighted-form-two}. At this
stage their $p,q$ are real numbers; they are not assumed to lift to $K$.
The existence of group-valued parameters is what must be proved.

For each form compute the fifteen coefficient differences of its ordered
list. Both sides have the same fifteen distinct coefficient pairs. These
can be checked by subtracting entries of the displayed lists; as real
expressions the two common multisets are
\begin{align*}
 \mathcal D_1=\{&q-3p,2q-3p,q-2p,2q-2p,3q-2p,\\
 &q-p,2q-p,3q-p,q,2q,p,q+p,2q+p,2p,q+2p\},\\
 \mathcal D_2=\{&2q-3p,2q-2p,3q-2p,q-p,2q-p,3q-p,\\
 &q,2q,p,q+p,2q+p,2p-q,2p,q+2p,3p-q\}.
\end{align*}
Put $t=q/p$. Two unequal coefficient forms $up+vq$ and $u'p+v'q$
can agree only when $v\ne v'$ and
$t=(u'-u)/(v-v')$. Enumerate these rational slopes for the $\binom{15}{2}$
pairs and retain the appropriate chamber. This yields precisely $3,4,5$
for the first form and $3/2,5/3$ for the second. Away from these slopes
there is exactly one matching, the formal coefficient matching. All
forward real differences are positive there, so their signs are $+1$.

At an exceptional $t=a/b$, evaluate each coefficient pair at the integer
parameters $(p,q)=(b,a)$ and group edges into equal-height bins. In a bin
of size $k$, enumerate all $k!$ bijections. A zero bin also has $2^k$
independent sign choices; a positive bin has only the positive sign.
Thus the exact finite input is:
\begin{center}\small
\begin{tabular}{lrrrrr}
\toprule
Form & Generic & First slope & Second slope & Third slope & Total\\
\midrule
$1$ & $1$ & $128$ at $3$ & $16$ at $4$ & $4$ at $5$ & $149$\\
$2$ & $1$ & $128$ at $3/2$ & $4$ at $5/3$ & --- & $133$\\
\bottomrule
\end{tabular}
\end{center}
These $282$ signed matchings exhaust all possible lifts of every
noncongruent projected pair, because only the five explicitly computed
slopes permit additional edge coincidences. This is an argument over all
real parameters, not a parameter sample.

For clarity, the complete \emph{acceptance predicate} for this finite
step is as follows. Reconstruct every matrix from its signed labels by
\eqref{eq:formal-matching-row}. For each of the seven cases, compare the
saved labelled matchings with the full independent bin-permutation or
edge-by-edge enumeration, including signs on zero edges; require exact
set equality and absence of duplicates. Each matrix has an assigned
representative with both integer identities
\eqref{eq:two-way-lattice-cert}. Each representative has a certified
diagonalization~\eqref{eq:smith-contract}. Finally, in the resulting
exact group coordinates require one of:
\begin{enumerate}
 \item two distinct labels on the same endpoint have identical coordinates;
 \item explicit $p,q$, endpoint interchange if used, translations $t_X,t_Y$
 and a sign $\eta$ satisfy
 \[
 \operatorname{multiset}(E_i)=t_X+\operatorname{multiset}(X(p,q)),\quad
 \operatorname{multiset}(F_i)=t_Y+\eta\operatorname{multiset}(Y(p,q)),
 \]
 with the sides interchanged when specified. These equalities are checked
 by evaluating all twelve coefficient formulas; they are not inferred
 from a successful or failed parameter search.
\end{enumerate}
The saved package passes this predicate. Its $21$ distinct labelled
integral lattices consist of $8$ binary $\Z^2$ presentations, $3$ binary
$\Z$ presentations, $8$ binary $\Z\oplus C_2$ presentations, all with
explicit Bloom identities, and $2$ $\Z$ presentations with forced
collisions. Distinctness of the representative lattices can also be
certified by a row whose class is nonzero in the other's certified quotient;
it is ancillary to coverage.

An actual signed matching of the group pair supplies $G_M\to K$.
The collision outcomes contradict the six-element hypothesis. Every
remaining outcome has a Bloom identity already in $G_M$; applying the
homomorphism gives that identity in $K$, with distinctness inherited from
$A,B$. This proves BF, with precisely the Iw and finite-predicate
dependencies stated above.
\end{proof}

For example the apparently new torsion boundary configuration
\[
 \{0,g,3g,3g+h,7g,8g+h\}/
 \{0,2g+h,4g+h,7g,7g+h,8g+h\},\qquad 2h=0,
\]
is already Bloom: take $p=g+h,q=3g$ and reflect $Y,X$, respectively,
about $8g+h$. Torsion in a matching presentation alone does not prove
that its images are non-shadows.

\begin{corollary}[Intrinsic free-quotient dichotomy]\label{cor:free-dichotomy}
For a finitely generated universal presentation, either the endpoint
multisets in its full free quotient are congruent, or BF applies and the
pair is group-valued Bloom. In the congruent case a single labelled
alignment is valid for every real solution.
\end{corollary}
\begin{proof}
If no fixed vector congruence holds in the free quotient, the equalities for each choice of
sign, permutation and anchor can hold after a real functional only
on a proper subspace of the dual. There are finitely many proposals, so
Lemma~\ref{lem:finite-subspaces} supplies a functional outside all of
them. Its two projected multisets are noncongruent and BF applies.
In the other case use Lemma~\ref{lem:fixed-alignment}.
\end{proof}

\section{The high-free-rank cover and its direct \texorpdfstring{L3$^*$}{L3*} form}
\label{sec:formal-high-rank}
\begin{lemma}[HR]\label{lem:hr}
Let $M$ be an original signed matching presentation whose universal
endpoints have six distinct points and are $\TI$ distinct. If $d\ge3$,
then, conditional on Iw and the finite certificate predicate below,
\[
 G_M\cong C_2\oplus\Z^3.
\]
After relabellings and independent rigid motions its two configurations
are exactly
\begin{equation}\label{eq:hr-master}
 C=\{0,p,q,p-q+h\},\quad
 A=C\sqcup\{r,r+h\},\quad
 B=C\sqcup\{p-r,p-r+h\},\quad2h=0,
\end{equation}
where $p,q,r$ form a basis of the free quotient and $h$ generates $C_2$.
Every admissible specialization of this family is directly L3$^*$.
There are no $\TI$-distinct six-set universal presentations of rank
$d=4$ or $5$.
\end{lemma}
\begin{proof}[Reduction and exhaustive finite contract]
For each $v\in V_M$, Iw says that its real multisets are either congruent
or a labelled Bloom pair. There are finitely many label permutations,
orientations and choices of origin. Each fixed Bloom choice defines a
linear subspace of the anchored ten-coordinate space of dimension at
most two: all coordinates are fixed integer-linear combinations of
its two parameters, and anchoring removes the translations. Thus $V_M$
is covered by finitely many intersections with congruence subspaces and
Bloom subspaces. Since $\dim V_M\ge3$, none of the Bloom subspaces can
contain it. Lemma~\ref{lem:finite-subspaces} gives one fixed congruence
alignment. Apply Lemma~\ref{lem:fixed-alignment} and retain the original
integral presentation; only its real heights have $b_i=a_i$.

By~\eqref{eq:height-codimension}, the rational height span $S$ has
rank at most two. The primitive directions of all nonzero rows
\eqref{eq:height-rows}, with an overall sign removed, form the following
explicit $S_6$-invariant set $\mathcal D$:
\begin{align*}
 &(1,-1,0,0,0,0) &&\text{and its $15$ projective permutations},\\
 &(2,-1,-1,0,0,0) &&\text{and its $60$ projective permutations},\\
 &(1,1,-1,-1,0,0) &&\text{and its $45$ projective permutations}.
\end{align*}
To verify that these are all possibilities, subtract two signed
incidence roots. Equal supports give zero or a multiple of a root;
one shared vertex gives a root or the second shape; disjoint supports
give the third shape. The counts are respectively $\binom62$,
$6\binom52$, and $\binom64\binom42/2$. Hence $|\mathcal D|=120$.
Dividing by a coefficient gcd here does not divide any row of $M$.

Enumerate the zero span, every one-direction span, and every two-direction
span of $\mathcal D$, identifying rational spans by exact row reduction.
This enumeration is exhaustive because a span of rank at most two has
a basis chosen from its generating rows. Partition these spans by all
$720$ permutations of the six coordinates. A simultaneous permutation
of the aligned endpoints and reanchoring induces an invertible integral
change of the original generators, so one representative of each orbit
suffices. The saved enumeration has $4176$ spans in $25$ orbits.

For each representative $S$, supply an integer height vector
$z=(0,z_1,\ldots,z_5)$ with the certificate property
\begin{equation}\label{eq:generic-height-contract}
 v\cdot z=0\quad\Longleftrightarrow\quad v\in S,
 \qquad v\in\mathcal D.
\end{equation}
Such a vector exists: each $v\notin S$ cuts a proper hyperplane in the
anchored rational kernel of $S$; Lemma~\ref{lem:finite-subspaces} over
$\mathbb Q$ gives a point outside all of them, and clearing denominators
makes it integral. The stored vector must be checked against all $120$
directions. Its purpose is to record all forced signed-edge equalities,
not to approximate an arbitrary real solution.

Enumerate every permutation between the fifteen edge occurrences of
$z$ and another copy of $z$ that preserves absolute edge differences.
For a positive-length bin the sign is uniquely determined by the two
oriented differences. For a zero bin both signs are included independently.
If a bin has size $k$, it contributes $k!$, with an additional factor
$2^k$ only at zero. Reconstruct the full integral matrix for every
matching by~\eqref{eq:formal-matching-row}. Every original presentation
in this branch occurs in this collection after the justified changes of
generators: each of its matching rows restricts to a row in $S$, so
\eqref{eq:generic-height-contract} makes that signed equality true at
$z$. Conversely any equality at $z$ represents the same rational forced
height relation. Extra presentations with different integral rank may
be enumerated, which is harmless; none is omitted.

The acceptance predicate for the finite cover has four parts:
\begin{enumerate}
 \item Independently reconstruct $\mathcal D$ and all rank-at-most-two
 spans; check that the $25$ orbit lists are disjoint and their union is
 the full $4176$-span set. Check~\eqref{eq:generic-height-contract}.
 \item For every orbit representative compare the entire saved labelled
 matching set with an independent edge-by-edge enumeration. Require
 exact set equality, no repeated record, and the bin-factorial count.
 \item Each discarded record must carry either $cM=e_i-e_j$ or
 $cM=f_i-f_j$ for distinct labels, or a permutation $\pi$, sign $\eta$
 and anchor $k$ with six identities
 \[
 c_iM=e_i-\eta f_{\pi(i)}+\eta f_k.
 \]
 Check that $\pi$ is a bijection and all coefficients are integers.
 These identities prove collision or rigid equivalence in every image.
 \item Every remaining record has matrices satisfying
 \eqref{eq:smith-contract} with diagonal
 $(1,1,1,1,1,1,2,0,0,0)$. In the exact quotient coordinates, evaluate
 a supplied origin, rigid alignment and $p,q,r,h$ and require all twelve
 master coordinates~\eqref{eq:hr-master} to equal the universal endpoints.
 Require $h\ne0$, $2h=0$, and determinant $\pm1$ for the $3\times3$
 free-coordinate matrix of $p,q,r$.
\end{enumerate}
The last determinant and the order of $h$ prove that these parameters
generate the whole $C_2\oplus\Z^3$, rather than a proper sublattice:
their free components give a basis, and $h$ generates the remaining
torsion. Thus the record certifies the universal family itself.

The saved files pass this predicate, with $14475$ total matchings:
$13916$ collision records, $555$ universal rigid-equivalence records,
and four master-family records. The discard records contain $17246$
integer row identities. All four survivors have the same height orbit,
with span generated by
\[
 z_0-z_1-z_2+z_5=0,\qquad z_3-z_4=0;
\]
the generic stored vector is $(0,8,1,3,3,9)$ and that orbit has $128$
matchings. Independent replay checks every record; the numerical totals
alone are not the acceptance predicate. The preceding exhaustiveness
argument and these identities prove the structural part of HR. Since
the survivors have free rank three, none has rank four or five.
\end{proof}

\begin{proof}[Direct algebraic L3$^*$ identity]
In the group ring of any abelian target with $2h=0$, put
$\mathsf P=x^p$, $\mathsf Q=x^q$, $\mathsf R=x^r$, $\mathsf H=x^h$.
Then $\mathsf H^{-1}=\mathsf H$ and
\[
 C=1+\mathsf P+\mathsf Q+\mathsf P\mathsf Q^{-1}\mathsf H,
 \qquad D=\mathsf R(1+\mathsf H).
\]
Direct expansion gives
\begin{equation}\label{eq:hr-cosymmetric-defect}
 C-\mathsf PC^*=(1-\mathsf H)
                    (\mathsf Q-\mathsf P\mathsf Q^{-1}),
 \qquad (C-\mathsf PC^*)(1+\mathsf H)=0.
\end{equation}
Take the fixed block to be $C$, the moving dyad to be $D$, and
$s=p-2r$. Its cross distribution is $E=CD^*$, so
\begin{align*}
 x^{-s}E
 &=\mathsf P^{-1}\mathsf R^2 C\mathsf R^{-1}(1+\mathsf H)\\
 &=\mathsf R\mathsf P^{-1}C(1+\mathsf H)
  =\mathsf R C^*(1+\mathsf H)=E^*.
\end{align*}
This is exactly the L3$^*$ hypothesis. The shifted dyad is
\[
 x^sD=\mathsf P\mathsf R^{-1}(1+\mathsf H),
 \quad\text{with points }\{p-r,p-r+h\},
\]
so its endpoint is precisely $B$. Six-point distinctness supplies the
required disjointness. This proves homometry and the named direct move
without a global reflection or a choice of half of $s$.

For comparison the earlier L5 description is also valid: globally reflect
$B$ by $t\mapsto p-t$ to obtain $D+\mathsf PC^*$, with $D$ fixed and
$C$ moving. Equation~\eqref{eq:hr-cosymmetric-defect} gives
$DC^*=\mathsf P^{-1}DC$, the strict L5 first-cross-term condition.
The direct L3$^*$ description is the shorter form of the same family.
In a cyclic target six-point distinctness forces $h\ne0$, hence the
modulus is even and $h=n/2$. Its free parameters may otherwise specialize
arbitrarily, subject to endpoint distinctness and $\TI$ distinction.
\end{proof}

The BF and HR finite packages are respectively
\code{results/2026-09-30-six-bloom-cylinders.json} and
\code{results/2026-09-30-six-free-rank/}. The checker
\code{tests/test_six_cylinder_branches_review.py} independently rebuilds
the signed matching sets and verifies the integer identities, without
importing the builders or Smith/Hermite routines. The separate Iw sources,
formulas and solver evidence remain
\code{notes/2026-09-30-six-integer-theorem.md} and its logged adversarial
review. The present exposition does not claim new novelty or formal
proof-assistant verification.

% BEGIN LOW-RANK INTEGRATION
\subsection{Congruent real projections at low free rank}
\label{sec:lr-expanded}

We now specify the finite verification used when the original signed
matching presentation has free rank one or two. All lattices in this
subsection are integral lattices in the same labelled $\Z^{10}$.
Only the auxiliary height arrangement is a rational arrangement.

\begin{lemma}[LR]\label{lem:lr}
Let $M$ be a full signed matching matrix and let
$d=10-\rank_{\mathbb Q}M\in\{1,2\}$. Suppose that one fixed relabelling,
reflection and reanchoring makes every real solution satisfy labelled
$B=A$. Every cyclic realization having six distinct points on each side
and different $\TI$ classes is a specialization of a Bloom, D, L2,
L3*, L4, L5 or L7 construction.
\end{lemma}

The proof consists of the mathematical reductions below and the stated
finite acceptance predicates. The counts describe accepted existing
certificate packages; they are not assumptions about a numerical trend.

\paragraph{The height direction arrangement.}
Put $A=(0,a_1,\ldots,a_5)$ after the fixed normalization of $B$.
Before anchoring, let $\alpha_{ij}=e_j-e_i\in\Z^6$, $i\ne j$.
Every substituted matching row is a difference of two such oriented
edges. Its nonzero primitive direction, considered up to sign, is a
permutation of
\[
 (1,-1,0,0,0,0),\qquad (2,-1,-1,0,0,0),\qquad
 (1,1,-1,-1,0,0).
\]
There are respectively $15$, $60$, and $45$ directions. Denote their
union by $\mathcal D$. Let $S$ be the rational span of the substituted
rows. The space of their height solutions is $S^\perp$, modulo the
constant vector. The projection from the original real solution space
to the anchored $A$ coordinates is an isomorphism: it is injective
because $B=A$, and substitution gives its inverse. Consequently
$\dim S=5-d$, hence $\dim S$ is three or four. Dividing out integer
content here describes $S$; it does not change any matching lattice.

\begin{lemma}[Exhaustive height representatives]
Every rank-$r$ span generated by $\mathcal D$ occurs, up to a vertex
permutation, by adjoining a direction to a representative of rank $r-1$.
For every such span $S$ there is an integral vector $h$, with $h_0=0$,
such that $v\cdot h=0$ for $v\in\mathcal D$ exactly when $v\in S$.
\end{lemma}
\begin{proof}
Choose a basis of $S$ from its generating directions and remove its last
member. A permutation taking the remaining span to a representative
also takes the removed member into $\mathcal D$, since $\mathcal D$ is
permutation invariant. Induction proves the first assertion, starting
at the zero span. For the second, each $v\notin S$ defines a proper
rational hyperplane in $S^\perp$. A finite union of proper hyperplanes
cannot cover this vector space. Choose a rational vector outside their
union, subtract its first coordinate from every coordinate, and clear
denominators. All directions have coordinate sum zero, so anchoring
does not alter the required evaluations.
\end{proof}

The finite height check starts with all spans of ranks zero, one and
two, and verifies the extension induction by exact rational arithmetic.
For each saved representative it applies all $720$ vertex permutations,
checks orbit size and disjointness, and checks every extension of every
preceding representative. Every saved orbit must meet that extension
list. Independently, a span can be represented by its primitive exterior
minors: two full-rank row lists represent the same rational span exactly
when their exterior coordinate vectors are proportional. The accepted
rank-three table has $43\,770$ spans in $104$ orbits; the rank-four table
has $116\,401$ spans in $211$ orbits. Each saved $h$ is checked against
all $120$ directions using exact dot products and exact span membership.
Thus equal heights and equal signed edge heights are precisely the
coincidences forced by the representative span, without sampling error.

\paragraph{Roots retain the original integral relations.}
For one fixed $h$, orient every nonzero-height edge upward and group
these edges by their positive height differences. If a bin has $m$
edges, its compatible bijections are the $m!$ permutations. A full
matching respecting $h$ restricts to one element of the product of
these bin-permutation sets. Equal-height edges are left for completion.
Relabellings inside each equal-height fibre act independently on the
source and target lists. They preserve the completion problem and
are integral invertible changes of anchored generators. In particular,
when the zero-labelled vertex moves, reanchoring subtracts its new
coordinate; the inverse relabelling gives an integral inverse.

For an upward source edge $(i,j)$ and target edge $(k,l)$ the root row
is $\alpha_{ij}^A-\alpha_{kl}^B$ in anchored coordinates. Upward order
need not have $i<j$. To use the fixed $i<j$ convention, multiply the
whole row by $-1$ if the source edge is reversed and transfer that
sign to the target; reverse the target edge as necessary and change
its sign accordingly. These operations leave its integral row lattice
unchanged. No equation $v\cdot h=0$ is added simply because it holds.

The root acceptance check reconstructs all bins and all within-fibre
vertex permutations. For each saved bin permutation $p$, it constructs
its complete orbit $\{\beta p\alpha\}$ under the independent actions,
checks the recorded orbit size, and checks disjointness from earlier
orbits. Their union must have cardinality $\prod m!$. It reconstructs
every saved root matrix from its edge permutation, rather than trusting
the saved matrix. Across the $315$ height strata this accounts for
$4\,822\,240$ cross-edge bijections in $328\,473$ root orbits.

\paragraph{Exact group coordinates and cache equality.}
A node has a row matrix $H\in\Z^{k\times10}$ and represents
$G_H=\Z^{10}/\row_\Z(H)$. Its Smith certificate consists of
$U\in\Z^{k\times k}$ and $V\in\Z^{10\times10}$ with
\[
 \det U,\det V\in\{1,-1\},\qquad UHV=D.
\]
Here $D$ is a rectangular diagonal matrix with $r$ nonzero entries,
where $r=\rank H$. Thus $U$ is not necessarily $15\times15$:
its dimension changes with the node's actual row count. Keeping the
coordinates of rows of $V$ at indices $i<r$ with $|D_{ii}|>1$, reduced
modulo $|D_{ii}|$, and all indices $i\ge r$ gives the ten generator
coordinates in $T\oplus\Z^{10-r}$. Entries $|D_{ii}|=1$ contribute
no coordinate. The two anchors are zero. The verifier checks the
matrix identity, dimensions, determinants, rank, diagonal data and all
twelve recorded point coordinates by integer arithmetic.

The producer caches the row-Hermite matrix of the relation lattice.
This cache is local to one fixed height stratum and the same labelled
generators; neither rational span equality nor an unrecorded abstract
group isomorphism permits a merge. For every root normalization and
every parent-plus-branch normalization, the package supplies integral
matrices $C,C'$ satisfying
\[
 C H_{\rm source}=H_{\rm target},\qquad
 C'H_{\rm target}=H_{\rm source}.
\]
These equalities prove equality of integral row lattices in both
directions, including torsion, without trusting a Hermite routine.
Canonical minimality is useful for efficiency, not for the coverage
proof: two-way lattice equality is the required mathematical contract.

\paragraph{Paired cancellation and exhaustive branching.}
In $G_H$ form a bin for each unsigned class $\{z,-z\}$ of the fifteen
edge differences, retaining its labelled edge occurrences. Give the
class the deterministic representative $\min(z,-z)$ in Smith coordinates.
Cancel the minimum of the two multiplicities for each common class.
This cancels equally many occurrences on the two sides, including
order-two classes; it is not cancellation of vertices or of real heights.

If two points on one side coincide in $G_H$, discard the node as a
forced collision. If $B=s+\epsilon A$ in $G_H$, with
$\epsilon\in\{1,-1\}$, discard it as forced $\TI$. Both conditions
persist in every quotient. Otherwise an empty residual list is a
homometric terminal. At a nonterminal choose one residual $A$ edge
$(i,j)$. For each residual $B$ class use its first labelled edge
$(k,l)$ and include each sign $\epsilon$ satisfying
\[
 h_j-h_i=\epsilon(h_l-h_k).
\]
The child adjoins the actual integral row
$\alpha_{ij}^A-\epsilon\alpha_{kl}^B$. Both signs are included if
both satisfy this equality, as happens for zero heights.

\begin{lemma}[Coverage and termination of completion]
Suppose a full presentation $G_M$ in this stratum is a quotient of a
node $G_H$. Unless its realization has a collision or is $\TI$,
it factors through one recorded child of every nonterminal encountered.
Every path terminates after at most fifteen uncancelled-edge decreases.
\end{lemma}
\begin{proof}
The cancelled classes remain equal in $G_M$. Subtracting the same
multiset from the two equal full edge multisets therefore leaves equal
residual multisets in $G_M$. The chosen $A$ occurrence equals some
residual $B$ occurrence up to sign there. Replacing that occurrence by
the saved representative of its unsigned class changes at most the
sign, since equality up to sign already holds in $G_H$. Both possible
signs are tested. The common real height homomorphism is defined on
$G_M$, so the resulting signed height equality is among the branches.
The newly adjoined row holds in $G_M$ and gives the required factorization.

Existing unsigned equalities persist on passing to a quotient. The new
equality pairs one occurrence from each residual side that had not been
paired before. Even if further classes merge, a matching of occurrences
extending all previous cancellations and this new pair exists. Hence the
maximum cancellable multiplicity increases by at least one, and the
residual multiplicity strictly decreases. It starts at most fifteen.
Two-way lattice certificates make cache reuse a reuse of the same
labelled quotient, so this decrease also precludes cycles after merging.
\end{proof}

This argument uses the original presentation and its real homomorphism;
it does not posit a nonzero homomorphism from the finite cyclic target
back to $\R$. The root-orbit reduction may first replace $M$ by its
within-fibre relabelling; that presentation has the same realization
up to the recorded changes of labels and anchors.

Here is a complete abstract version of the completion checker:
\begin{verbatim}
for each height stratum:
  verify height orbit cover, generic heights, and cross-root orbit cover
  for each root: check both integral inclusions into its saved node
  visit each reachable node once:
    verify Smith coordinates and all point/edge classes
    recompute residual lists by paired multiplicity cancellation
    if collision: verify two distinct labels have equal coordinates
    elif T/I: verify the saved sign and shift give the entire other set
    elif homometric: require empty residuals and six points on both sides
    else:
      require chosen edge belongs to the residual A list
      regenerate every height-compatible (B representative, sign)
      require the saved branch list equals this list, without duplicates
      for each branch: check its exact row and both lattice inclusions
      require the child's residual multiplicity is strictly smaller
  require every saved node is reachable and every root orbit is represented
\end{verbatim}
The package must report \code{complete=true}, no error, and certified
normalizations; a capped or interrupted search is rejected. The accepted
cover has $10\,602$ nodes: $7\,429$ forced $\TI$, $2\,207$ collisions,
$620$ homometric terminals and $346$ branching nodes, with $974$ branch
edges. The $658\,894$ integral inclusion identities are exactly twice
the number of roots plus branch edges.

\paragraph{Terminal mechanism predicates.}
For every homometric terminal the mechanism atlas must contain exactly
one record with the same stratum, node identifier, points and group orders.
Write $\widetilde B=s+\epsilon B$ for its certified alignment. For a
block certificate check a nonempty disjoint partition $A=U\sqcup W$,
its proposed moved block $W'$, and $\widetilde B=U\sqcup W'$.
With $P=UW^*$, the exact integer coefficient tests are
\[
\begin{array}{c|c|c}
\text{label}&W'&\text{acceptance identity}\\\hline
\mathrm{L2}&x^tW&x^tP=P\\
\mathrm{L3*}&x^tW&x^{-t}P=P^*\\
\mathrm{L4}&x^tW&2t=0,\quad x^t(P+P^*)=P+P^*\\
\mathrm{L5}&x^tW^*&P=x^{-t}UW .
\end{array}
\]
The group-ring soundness proofs given above apply literally. A verifier
expands these distributions as integer counters; it never accepts a label
merely because an endpoint autocorrelation check succeeds.

A Bloom certificate supplies $p,q$, anchor and alignment. Evaluate the
six displayed Bloom formulas and require exact endpoint equality with
multiplicity. A D certificate supplies an anchor, four-set $C$ and steps
$a,b$. Check the two endpoint formulas and expand
\[
 K=C+x^{a+b}C^*+(1+x^a)(1+x^b),\qquad
 (1-x^a)(1-x^b)K=0.
\]
The cyclic periodic-kernel lemma proved above converts this exact identity
to the stated nonnegative periodic-weight D grammar. If specialization
makes the two equal-sum dyads share a point, they share both points and
the endpoints agree; otherwise the four corners remain distinct.

An L7 certificate supplies a generator $h$ of a subgroup $H$ of order
$m\in\{2,4,6,12\}$. Enumerate its $m$ multiples and check their
distinctness and $mh=0$. Form $S=A+H$, require $|S|=12$, check exactly
$m/2$ points of $A$ in every occupied $H$-coset, and require
$\widetilde B=S\setminus A$. Its previously proved specialization
lemma is essential: a surviving non-$\TI$ six-point cyclic image
preserves the subgroup order and occupied cosets. Thus the image is a
literal L7 construction, not only a transported polynomial identity.

The producer's accepted $611$ direct terminal certificates comprise
$271$ L4, $135$ L3*, $114$ L7, $53$ L5, $16$ L2, $14$ Bloom and $8$ D.
A different independent search priority supplies $481$ D, $71$ L5,
$24$ L3*, $20$ L4, $14$ Bloom and one L2, again leaving exactly nine
terminals. These distributions are not invariants or a disjoint
classification; acceptance depends on the supplied identities.

\paragraph{All cyclic images of noncyclic torsion.}
For any remaining terminal write $G_H=T\oplus\Z^r$ with
$T=\bigoplus_{i=1}^k C_{d_i}$ and $E=\operatorname{lcm}(d_i)$.
Enumerate all tuples $j_i\in\{0,\ldots,d_i-1\}$, set
$c_i=(E/d_i)j_i$, $g=\gcd(E,c_1,\ldots,c_k)$ and $q=E/g$.
The character has actual image order $q$, and its surjective image map is
\[
 \theta_j(t_1,\ldots,t_k)=
 \frac{\sum_i c_it_i}{g}\pmod q.
\]
Retain the free coordinates unchanged. If $q=1$ omit the torsion coordinate.

\begin{lemma}[Actual-image character factorization]
Every map $f:T\oplus\Z^r\to C_n$ factors through one of the displayed
quotients $C_q\oplus\Z^r$, with $q\mid n$.
\end{lemma}
\begin{proof}
The image $f(T)$ is cyclic of order $q$ dividing both $E$ and $n$.
Identify it with $C_q$, and embed $C_q$ in $C_E$ by multiplying by
$E/q$. Each generator of order $d_i$ then maps to an element killed by
$d_i$, hence to $(E/d_i)j_i$ for one enumerated $j_i$. Dividing the
resulting character by its common divisor $g=E/q$ recovers the chosen
map onto $C_q$. Extend its factor map by the original images of the free
generators. These images are unrestricted and need not be disjoint from
the torsion image. Requiring $C_E$ itself to embed into $C_n$ would be
incorrect; only the actual image $C_q$ is used.
\end{proof}

The checker requires every character tuple exactly once, recomputes
$c_i,E,g,q$ and all projected point coordinates, then verifies a collision,
a $\TI$ witness or one of the preceding mechanism certificates.
Seven residual terminals have $T=C_2\oplus C_2$: each of their four
characters forces a collision. Two have $T=C_2\oplus C_8$: each has
sixteen characters, eight colliding and eight admitting L7 certificates
(independently also D). Thus all $60$ quotients give $44$ collisions and
$16$ mechanisms. Coverage, specialization and this factorization prove LR.

\subsection{The bounded finite-source branch}
\begin{lemma}[FT]\label{lem:ft}
For every $6\le q\le135$, each homometry family of six-subsets of
$C_q$ is connected by the stated grammar.
\end{lemma}
\begin{proof}[Proof with its finite acceptance contract]
Enumerate six-subsets in two independent ways. First enumerate every
positive gap word $(g_0,\ldots,g_5)$ with sum $q$, keeping precisely the
lexicographically least word among its rotations and reversals. Equivalently
require $g_0$ minimal and then test all twelve gap images. Cumulative sums
give its anchored point tuple. Gap rotations and reversals are exactly
translations and reflections of the cyclic point set, so this visits every
$\TI$ class once, including symmetric classes. The second traversal
enumerates every $0<a_1<\cdots<a_5<q$ and keeps the least sorted tuple
among all images anchored at one of its six points, with either sign.
Any rigid image containing zero has such an anchor, so this test is complete.
Lexicographic comparison of point tuples and their gap words first differs
at the same gap, so the two canonical conventions agree.

For the first traversal compute the fifteen distances
$\min(a_j-a_i,q-a_j+a_i)$. For the second compute directed membership
counts $\#\{x\in A:x+d\in A\}$ for $1\le d\le q/2$, halving the
antipodal entry. Expand these counts into the sorted fifteen-distance list.
Both lists are exact interval inventories. Group complete inventories,
retaining groups of at least two classes. The implementation packs fifteen
distance bytes in two words and five point bytes in another; for
$q\le135$ these encodings are injective. Equality uses both full distance
words, not a hash. Two traversals must agree family-for-family and in
total class count at every modulus, not merely in aggregate totals.

An additional class-count check uses Burnside's formula. A rotation by
$t$ has $c=\gcd(q,t)$ cycles of length $\ell=q/c$ and fixes
$\binom{c}{6/\ell}$ six-subsets when $\ell\mid6$, otherwise zero.
A reflection with $f$ fixed points fixes
$\sum_{j:(6-j)\text{ even}}\binom f j
\binom{(q-f)/2}{(6-j)/2}$, where $f=1$ for odd $q$, and $f=2$ or
$0$ according to even or odd reflection parameter for even $q$.
Sum over all $2q$ group elements and divide by $2q$.

Each census checkpoint records \code{methods}, \code{summary},
\code{independent}, \code{families} and \code{source_sha256}.
Each family stores its complete \code{members} and \code{icv}.
Reference regressions check the immutable implementation through $22$;
additional controls attack the $31$--$33$ and $63$--$65$ boundaries.

For every unordered pair within every family, the mechanism file stores
either a direct certificate or an explicit gap entry. Certified pairs and
gaps must partition the complete pair set, without duplicates or omissions.
A direct record stores canonical \code{a}, \code{b} and \code{label}.
Bloom records store \code{parameters} $[p,q']$; R records store a seed
index and multiplier, with its seed-order relation checked. L7 records
store subgroup order \code{h}; unit records store \code{u}, with
$\gcd(u,q)=1$, endpoint equality and invariant autocorrelation checked.
Block/D records store \code{certificate} $[\lambda,\epsilon,s,\mu,t]$:
$\mu$ is a six-bit mask selecting $U$ from the recorded $A$, while
$\epsilon,s$ align $B$. Reconstruct both blocks and verify the preceding
coefficient identities. The older \code{swap} label additionally requires
its halved-shift solvability condition; it is therefore contained in L3*.
For D reconstruct its equal-sum dyads and check the displayed mixed
difference of its twelve-unit weight distribution. The cyclic kernel lemma
then supplies the required nonnegative periodic decomposition.

Finally, form an undirected graph on each family's members with its
verified direct certificates as edges. Require one connected component,
using an independent union-find or exhaustive reachability calculation.
For every gap pair save and verify an explicit path of these edges; no
direct construction is required for that pair. The accepted tables have
$725\,132$ families and $728\,424$ unordered pair edges, $728\,351$
direct certificates and $73$ gaps, all connected by verified paths.
Moduli $6$--$11$ have no nontrivial family. These exact finite conditions,
together with soundness of the generators, prove FT.
\end{proof}

\begin{lemma}[Inflation]\label{lem:inflate}
If $q\mid n$, any generating path of six-subsets of $C_q$ embeds to
a generating path in $C_n$ under $x\mapsto(n/q)x$. Distinct $\TI$
classes stay distinct. Consequently FT and the rank-ten finite-source
bound settle the original free-rank-zero branch at every modulus.
\end{lemma}
\begin{proof}
The embedding is injective. It preserves every point distinction,
partition, subgroup, periodic-weight distribution and group-ring identity,
so it transports Bloom, D, block, coset and rigid-template edges. For a
unit $u$ modulo $q$, choose $u'$ congruent to $u$ modulo $q$ and congruent
to $1$ modulo every prime dividing $n$ but not $q$. The moduli in these
extra congruences are coprime to $q$, so CRT supplies $u'$. Primes dividing
$q$ cannot divide $u'$ because they do not divide $u$; the other primes
were excluded explicitly. Thus $u'$ is a unit modulo $n$ and acts on
the embedded subgroup exactly as $u$. Autocorrelation is zero off that
subgroup, hence the multiplier condition also transports.

If two nonempty subsets of the embedded subgroup $H$ satisfy
$B=s+\epsilon A$ in $C_n$, choose one point on each side in this equality.
Their difference shows $s\in H$. The same equality is therefore a rigid
equivalence inside $H\cong C_q$, proving preservation of vertices.
For an original rank-ten presentation $G_M$ the group is finite, with
$|G_M|\le135$ by the minor bound. Its cyclic image has order
$6\le q\le135$ and divides $n$. Independently anchored inputs lie
inside this subgroup, so FT applies and its path embeds as just proved.
Restore the original anchor translations at the endpoints.
\end{proof}

\paragraph{Reproducibility and review boundary.}
The height tables are in \code{results/2026-09-30-six-free-rank/};
the compressed cross/root/node certificates are in
\code{results/2026-09-30-six-free-dag-all/}; the terminal atlas is
\code{results/2026-09-30-six-free-mechanisms/atlas.json}.
The structural checker is \code{tests/test_six_free_dag_review.py};
the independent terminal checker is
\code{tests/test_six_low_rank_mechanisms_review.py}.
They check different obligations and both are required. The additional
height checker is \code{tests/test_six_free_rank_growth.py}.
Finite census and edge checkpoints are respectively in
\code{results/2026-09-30-six-large-census/} and
\code{results/2026-09-30-six-pair-mechanisms/}; the finite bridge audit is
\code{tests/test_six_finite_torsion_review.py}.
The supplied independent review of the earlier draft rechecked height
counts and consequences but explicitly did not replay the internal LR DAG.
Its review must not be cited as such a replay. The separately implemented
in-house structural and terminal checks are the certificate evidence for
that branch. These finite computational claims remain part of a
computer-assisted proof, not a proof-assistant-kernel or novelty claim.

% END LOW-RANK INTEGRATION
\section{Completeness of the generating grammar}
\begin{proof}[Proof of Theorem~\ref{thm:main}]
Soundness is Proposition~\ref{prop:sound}. $\TI$-equivalent inputs have
a length-zero path. For the converse on distinct vertices choose any
compatible universal signed matching for a $\TI$-distinct homometric input.
If $d=0$, Corollary~\ref{cor:finite} and FT give a finite source path;
Lemma~\ref{lem:inflate} lifts it. If $d>0$ and some real projection is
noncongruent, BF gives B directly. Otherwise one fixed congruence holds
on the whole real solution space, as in Section~\ref{sec:universal}.
HR handles $d=3,4,5$, while LR handles $d=1,2$. Specialization and
Lemma~\ref{lem:l7} retain the named admissible generators. These cases
exhaust the original presentation dimensions and both projection branches.
\end{proof}

\begin{definition}[Integer shadow]
A cyclic Z-pair is an integer shadow if it is the reduction of a
homometric pair of six-element integer sets. It is purely cyclic if
no such lift exists. A cyclic construction label alone does not determine
this distinction.
\end{definition}

\section{Integer shadows and the large-prime boundary}\label{sec:shadowfull}

\begin{proposition}\label{prop:shadowexact}
Assuming the line obligation U, a cyclic Z-pair is an integer shadow
if and only if representatives are an admissible modular Bloom pair.
\end{proposition}
\begin{proof}
Apply the distinct-point line theorem I to any integer lifts and reduce
its integer rigid motions and recovered integer parameters modulo $n$.
For the converse lift the two modular parameters to any integers $p,q$.
Six distinct residues in each formula force six distinct integer terms.
The formal Laurent identity for B is valid in $\Z[x,x^{-1}]$, so those
integer sets are homometric. An integer congruence would reduce to a
cyclic $\TI$ congruence; the latter has been excluded. This also explains
why exhaustive testing of the $n^2$ modular parameter pairs, with rigid
alignments, is an exact shadow test despite unbounded integer lifts.
\end{proof}

\begin{corollary}\label{cor:largeprime}
If every prime divisor of $n$ is greater than 131, every cyclic six-point
Z-pair is an admissible modular Bloom image, under the same line-solver
qualification.
\end{corollary}
\begin{proof}
Take its universal presentation $G_M=\Z^d\oplus T$ and map
$f\colon G_M\to C_n$. The torsion image has order dividing both $n$
and $|T|\le135$. The smallest prime above 131 is 137, so this image
is trivial. The pair factors through the free quotient. Its six vertices
in each list remain distinct, since their images in $C_n$ are distinct.
Here $d>0$, because a finite presentation with trivial image cannot
contain a six-set.

Choose a basis of the free quotient and any integer representatives
$c_1,\ldots,c_d$ of its images under $f$. The integer functional
$\lambda(w)=\sum_j c_jw_j$ reduces modulo $n$ to $f$. Any within-endpoint
difference $w$ has nonzero image under $f$, so $\lambda(w)$ is a nonzero
integer. Thus $\lambda$ is injective on each six-element endpoint, with
no genericity assumption. Homomorphic transport gives integer homometric
six-sets reducing to the original cyclic pair.
Proposition~\ref{prop:shadowexact} finishes.
\end{proof}

This sufficient boundary is not an enumeration formula. In particular,
deducing $(p-1)(p-11)/12$ at primes requires a separate proof of parameter
admissibility, all identifications between parameter images, and all
exceptional stabilizers. That counting problem is not settled merely by
the large-prime corollary or by matching finite numerical values.

\subsection{What the rigid-template certificates additionally prove}
The soundness of an R edge requires only equality of its fifteen unsigned
differences modulo its listed $q$. The stronger claim that a faithful seed
image is purely cyclic uses an additional, pattern-restricted statement.
We give its exact finite contract to keep these two uses separate.

For each row $(q,X,Y)$ of Table~\ref{tab:rigid}, order its six coordinates
as printed and exhaust all bijections of $\mathcal E$ and all signs
compatible with those coordinates modulo $q$. Compatibility means the
row relation~\eqref{eq:formal-matching-row} evaluates to zero modulo $q$;
both signs are retained for an involution. The saved arithmetic witnesses
must, for every such matching and with no omitted or duplicate label list,
supply integer unimodular $U,V$ satisfying
\[
 UMV=\operatorname{diag}(1,1,1,1,1,1,1,1,1,q),
\]
with the five additional rows zero. If $w$ is the last column of $V$ and
$v_{\rm seed}$ is the ten-coordinate anchored seed vector, also require
\[
 \gcd(w_1,q)=1,\qquad w_i\equiv w_1(v_{\rm seed})_i\pmod q
 \quad(1\le i\le10).
\]
The seed's first nonanchored coordinate is one. These predicates are
checked by integer multiplication, determinants and congruences. The
numbers of compatible matchings, in table order, are
$192,96,48,48,48,48,16,12,8,8,4,4,4$, totaling $536$.

\begin{proposition}[Rigidity restricted to the seed matching patterns]
Under the preceding finite verification, every realization of a compatible
seed matching in any abelian group is the displayed seed multiplied by
one element $t$ with $qt=0$. Every faithful cyclic image of each listed
seed is purely cyclic.
\end{proposition}
\begin{proof}
The unimodular diagonal identity identifies its universal presentation
with $C_q$. All old vertices are multiples of the generator, with
coefficients $w_i$. Replacing the generator by $w_1$ times itself is
invertible modulo $q$ and gives exactly $v_{\rm seed}$. Thus every
realization has the asserted form. The allowed orders in the table are
checked by reducing both seed lists to each divisor of $q$, requiring
six distinct coordinates and different $\TI$ classes. By the subgroup
equivalence argument in Lemma~\ref{lem:inflate}, that check is valid in
every ambient cyclic group.

If the seed itself had homometric six-point integer lifts, their equal
absolute edge lengths would give one of the exhausted signed matchings
after reduction modulo $q$. Its anchored integer coordinates would lie
in the kernel of a rank-ten $M$, hence would all be zero, contradicting
six distinct points. Now let a faithful image lie in $C_n$, with
$t=(n/q)u$ and $\gcd(u,q)=1$. Anchor hypothetical integer lifts at their
points mapping to zero. All their coordinates are divisible by $n/q$.
Divide by $n/q$, then multiply by any integer inverse of $u$ modulo $q$.
These operations preserve integer homometry and distinctness and would
give integer lifts of the original seed, a contradiction.
\end{proof}

The proposition concerns the listed signed matching patterns, not all
six-point pairs. Proper images can admit additional matching patterns;
their pure cyclicity does not follow from the faithful-image argument.
In particular, this verifies the 21-EDO benchmark's pure cyclicity without
assuming that a cyclic factor flip is an integer-line factorization.

The 21-EDO pair
\[
\{0,1,3,7,10,15\}/\{0,1,4,7,14,16\}
\]
is the fifth rigid row and has separate exact non-shadow verification.
It also has a complement-conjugated swap, a cyclic integral factor flip
and a difference-set exchange explanation. Multiple explanations can
coexist. In particular, a factor flip in $\Z[C_n]$ need not be a
homometric factorization over the integer line.

Inflation is the injective subgroup embedding of Lemma~\ref{lem:inflate}.
General homomorphic specialization may impose additional torsion
relations. Compositions are paths through intermediate $\TI$ classes.
These notions are retained separately even where their parameter images
overlap.

\section{Certificate inventory and reproducibility}\label{sec:certificates}
All paths below are relative to the supplied project checkout. The companion
\code{evidence-manifest.json} binds the listed proofs, attacks, scripts and
saved certificate artifacts to SHA256 hashes. This draft has no public
archival DOI or immutable public artifact URL; those must precede a submission
claim of publicly reproducible computations. The review package contains
sources, proofs and dated six-point mathematical outputs. It also preserves
historical exploratory outputs; only the complete certificate packages named
by the proof are completeness dependencies.

\begin{center}\small
\begin{tabular}{ll}
\toprule
Component & Written proof / scope of independent evidence\\
\midrule
I/Iw & Integer and weighted real-line solver encodings\\
BF & 282 signed matchings; 21 integral presentations\\
HR & 4176 spans; 25 orbits; 14475 signed matchings\\
C/C2 & 3003 ten-edge graphs; two spanning-tree counts\\
R & 536 universal matching certificates; all allowed orders\\
FT & Two complete six-only enumerations through 135; exact paths\\
LR & 315 strata; full quotient cover and terminal arithmetic\\
G & Separate whole-argument attack and interface repairs\\
Driver & Portable certificates; separate implementation review\\
\bottomrule
\end{tabular}
\end{center}

The logical component documents have prefix
\code{notes/2026-09-30-six-} and respective suffixes
\code{integer-theorem.md}, \code{bloom-cylinders.md},
\code{free-rank.md}, \code{completeness.md}, \code{templates.md},
\code{finite-torsion.md}, \code{free-dag.md} and
\code{generation.md}. The synthesis's proof map pairs each with its
separate attack. The elementary D proof and cyclic kernel parameterization
are in \code{shadow-theorem.md} with \code{shadow-review.md}.

The LR structural replay checks 658894 integral lattice-inclusion identities,
all 974 branch edges, all initial matching orbits and completion at 315
expected strata. Smith witnesses check $UMV=D$ and unimodularity.
Two-way row identities certify normalization without trusting a library's
claimed normal form. Its terminal arithmetic checker imports no builder,
recognizer, Smith or Hermite routine. The independently chosen mechanism
priority gives a different noncanonical label distribution while still
covering every terminal.

For clarity, the stronger R rigidity verification is needed for the
faithful-image non-shadow statement, not for the soundness of an R edge.
The latter follows by checking the fifteen unsigned differences in its
printed seed row. The real-line obligation U, the spanning-tree bound,
the BF/HR finite covers, the full LR cover and the finite-source graph
cover are the actual completeness dependencies. Generator discovery
priorities, random controls and larger-modulus numerical trends are not
additional completeness assumptions.

Use the pinned Python 3.12/NumPy 2.5 environment. The default Python 3.14 /
NumPy 1.26 environment on the original host has a demonstrated array-mutation
defect and is not evidence for these computations. A focused replay is:
\begin{quote}\small\ttfamily
.venv/bin/python tests/test\_env.py\par
.venv/bin/python tests/run\_tests.py\par
.venv/bin/python tests/test\_six\_cylinder\_branches\_review.py\par
.venv/bin/python tests/test\_six\_free\_rank\_growth.py\par
.venv/bin/python tests/test\_six\_free\_dag\_review.py --resume --require-complete\par
.venv/bin/python tests/test\_six\_low\_rank\_mechanisms\_review.py\par
.venv/bin/python tests/test\_six\_finite\_torsion\_review.py
\end{quote}
The saved certificate driver also accepts arbitrary modulus. For example:
\begin{quote}\small\ttfamily
.venv/bin/python src/six\_generate.py 21\par
\hspace*{1em}0,1,3,7,10,15 0,1,4,7,14,16 --out results/example.json\par
.venv/bin/python src/six\_generate.py --replay results/example.json
\end{quote}
The first command is shown on two lines for typesetting and is entered as
one shell command. The driver is a reviewed implementation of G, not a
substitute for its completeness proof. Full regeneration instructions,
solver inputs and completion flags are in the component proofs. Long replay
jobs have resumable checkpoints bound to input hashes.

\section{Musical examples and limitations}\label{sec:trust}
The four-note teaching example has interval vector $(1,1,1,1,1,1)$.
A five-note example fixes $\{0,4,8\}$ while translating $\{3,5\}$ to
$\{5,7\}$; its six directed cross differences fill the odd residues once,
so L2 preserves them. The endpoints have vector $(2,1,2,3,2,0)$ and
different $\TI$ classes. Adjoining the same pitch to a Z-pair is not a
general construction: adding 8 to the displayed four-note pair destroys
homometry. Pair counts increase from 3 to 6 to 10 to 15, but this teaching
progression is not an induction on cardinality.

In $C_{12}$ all fifteen six-note Z-pairs are explained by classical
complementation. An ordinary-piano study and a separate retuned 21-EDO
miniature demonstrate how equal inventories can support contrasting
musical gestures. Registered integer Bloom passages use actual key positions,
not a claim that their reduction modulo twelve has six distinct pitch
classes. The narrated lesson and editable scores are explanatory companions,
not mathematical proof certificates.

The result's trust base includes exact integer arithmetic in audited C/Python
code, independently implemented traversals and the independently encoded
solver results for Iw. The separate synthesis reviewer refreshed solver
input/proof hashes but did not rerun every component or externally check the
solver's proof objects. There is no full Lean formalization. These limits
are material to the meaning of ``computer-assisted proof.''

Two online priority searches read relevant primary sources and located no
matching full arbitrary-modulus six-point generating theorem. Access gaps
include several potentially overlapping full papers and direct bibliographic
database searches. Absence from those searches is not evidence of originality.
This draft's theorem is the project's in-house result; attribution of any
new contribution awaits specialist literature review. Unread historical
sources are not used as axioms or cited as verified proofs here.

The next mathematical tasks are to remove redundant generators, describe
overlapping parameter images, seek useful unique forms and derive counting
formulas. Higher cardinalities remain outside this theorem. The project's
separate all-cardinality census and higher-deck definition issues have no
role in the proof and retain their recorded limitations.

\section*{Authorship, AI assistance and status}
The author, Hunter Bown, set the research questions and standards of this
programme and is responsible for its claims. The mathematics, code,
exposition and the separate in-house adversarial reviews were produced with
substantial assistance from AI language-model agents working under the
author's direction; the agents wrote first drafts of most proofs and of this
text. Those in-house reviews are not external peer review.

Two independent safeguards were added on 4 October 2026 and are part of the
accompanying repository. First, the soundness of L2, L3$^*$, L4, L5 and L7,
the exact dyad identity, the Bloom factorization and the
finite-union-of-subspaces step of Section~\ref{sec:grammar} onward are proved in
Lean~4 against Mathlib, with no \texttt{sorry}, in the directory
\code{formal/}. Second, the three cvc5 proof certificates behind Theorems~I and~Iw were
replayed in Ethos, a proof checker separate from cvc5, against the CPC
signature of the cvc5 source, with each proof's assumptions matched to the
saved problem (\code{tools/proof\_replay/}); an attempted solver-free
re-derivation by branch-and-bound over sorted interval orders succeeds for
four and five atoms but does not terminate for six
(\code{tools/iw\_enumeration/}). Four historical sources (Yovanof 1988, Soderberg 1995, Patterson 1944,
Bullough 1961/1964) were checked at first hand to the extent their texts
were accessible; Yovanof's thesis classifies only distinct-distance rulers
\cite{yovanof}, and none of the sources read states a six-point generating
theorem for all cyclic groups.

Before submission or public claims, a human specialist should still inspect
the universal presentation cover, the weighted-line encodings, the lattice
checkers, the specialization lemmas and the prior-art comparison. This
manuscript has not been submitted or published.

\begin{thebibliography}{99}
\small
\bibitem{rosenblatt}
J. Rosenblatt, \emph{Phase retrieval}, Communications in Mathematical
Physics \textbf{95} (1984), 317--343.
\href{https://doi.org/10.1007/BF01212402}{doi:10.1007/BF01212402}.
Four-point statement and attribution: Theorem 3.9 and the associated remark.

\bibitem{ej}
W. Q. Erickson and N. B. Jones, \emph{Homometric subsets of $\Z_n$ with
cardinality 5: classification and enumeration}, arXiv:2412.08997,
version 4, 27 August 2026.
\href{https://arxiv.org/abs/2412.08997v4}{arXiv:2412.08997v4}.
The five-point classification and final-section questions were read;
preprint status is distinguished from external review of the present result.

\bibitem{pg}
E. Postpischil and P. Gilbert, \emph{There Are No New Homometric Golomb
Ruler Pairs with 12 Marks or Less}, Experimental Mathematics
\textbf{3}(2) (1994), 147--152.
\href{https://doi.org/10.1080/10586458.1994.10504286}{\nolinkurl{doi:10.1080/10586458.1994.10504286}}.
Classical six-mark family and symbolic matching method; its Golomb-ruler
scope is not used to infer repeated-distance or cyclic completeness.

\bibitem{jj}
F. Jedrzejewski and T. Johnson, \emph{The structure of Z-related sets},
in \emph{Mathematics and Computation in Music}, LNCS \textbf{7937}
(2013), 128--137.
\href{https://doi.org/10.1007/978-3-642-39357-0_10}{doi:10.1007/978-3-642-39357-0\_10}.
Author version read:
\href{https://arxiv.org/abs/1304.6608}{arXiv:1304.6608}.
General constructions and finite enumeration.

\bibitem{mandereau}
J. Mandereau, D. Ghisi, E. Amiot, M. Andreatta and C. Agon,
\emph{Discrete phase retrieval in musical structures}, Journal of
Mathematics and Music \textbf{5}(2) (2011), 99--116.
\href{https://doi.org/10.1080/17459737.2011.608820}{doi:10.1080/17459737.2011.608820}.
Read through the author preprint; final typeset numbering was not assumed.

\bibitem{mandereaua}
J. Mandereau, D. Ghisi, E. Amiot, M. Andreatta and C. Agon,
\emph{Z-relation and homometry in musical distributions}, Journal of
Mathematics and Music \textbf{5}(2) (2011), 83--98.
\href{https://doi.org/10.1080/17459737.2011.608819}{doi:10.1080/17459737.2011.608819}.
Quotient transport and complementation were read through the author preprint.
\bibitem{yg}
G.~S. Yovanof and S.~W. Golomb,
\emph{The Polynomial Model in the Study of Counterexamples to
S. Piccard's Theorem}, Ars Combinatoria \textbf{48} (1998), 43--63.
Family N, p.~47; endpoint product labels are interchanged relative
to the displayed factors.

\bibitem{lss}
P. Lemke, S.~S. Skiena and W.~D. Smith,
\emph{Reconstructing Sets From Interpoint Distances},
DIMACS Technical Report 2002-37 (2002).
Published in \emph{Discrete and Computational Geometry},
Springer (2003), 597--631.
The primary report was read; Table~1 supplies the qualified
Hosemann--Bagchi historical attribution.

\bibitem{ch}
C. Callender and R. Hall,
\emph{Crystallography and the structure of Z-related sets},
Society for Music Theory annual meeting, Nashville, Tennessee,
7 November 2008, conference handout.
Cases~3--4; the Bullough attribution is read through this handout.

\bibitem{goyette}
J. Goyette,
\emph{The Z-Relation in Theory and Practice},
Ph.D. thesis, University of Rochester (2012).
Chapter~3, especially Formulas~3.2 and~3.5; the proposed criteria
are explicitly empirical (p.~122, footnote~2).

\bibitem{yovanof}
G.~S. Yovanof, \emph{Homometric Structures}, Ph.D. dissertation,
University of Southern California (1988).
\href{https://doi.org/10.25549/usctheses-c3-279207}{doi:10.25549/usctheses-c3-279207}.
Read in part in the library viewer (contents, abstract, \S\S2.7, 4.1, 4.3--4.4, 7.1);
Family~N, p.~156, is proved unique for six-mark rulers with all distances distinct.

\bibitem{soderberg}
S. Soderberg, \emph{Z-Related Sets as Dual Inversions}, Journal of Music
Theory \textbf{39}(1) (1995), 77--100.
\href{https://doi.org/10.2307/843899}{doi:10.2307/843899}.
Page~77 read; the construction is otherwise known here through Goyette.

\bibitem{patterson}
A.~L. Patterson, \emph{Ambiguities in the X-Ray Analysis of Crystal
Structures}, Physical Review \textbf{65} (1944), 195--201.
\href{https://doi.org/10.1103/PhysRev.65.195}{doi:10.1103/PhysRev.65.195}.
Abstract read; contents known through secondary sources.

\bibitem{bullough1}
R.~K. Bullough, \emph{On homometric sets. I. Some general theorems},
Acta Crystallographica \textbf{14} (1961), 257--268.
\href{https://doi.org/10.1107/S0365110X61000838}{doi:10.1107/S0365110X61000838}.
First page read; the remainder was inaccessible.

\end{thebibliography}
\end{document}
